\section{Einleitung}\label{einleitung}

\subsection{Zweck der
Testspezifikation}\label{zweck-der-testspezifikation}

Diese Testspezifikation beschreibt die Planung, Durchführung und
Bewertung der Tests für den im Rahmen der Projektarbeit entwickelten
Coasterbot-Prototyp.

Ziel der Tests ist es, die im Lastenheft definierten Anforderungen
systematisch zu überprüfen und nachzuweisen, dass die technische
Umsetzung gemäß dem Pflichtenheft funktioniert.

Die Testspezifikation stellt eine Verbindung zwischen folgenden
Dokumenten her:

\begin{itemize}
\tightlist
\item
  Anforderungen aus dem Lastenheft
\item
  technischen Umsetzungen aus dem Pflichtenheft
\item
  implementierten Systemkomponenten
\item
  durchgeführten Prüfungen
\end{itemize}

Dadurch wird sichergestellt, dass jede relevante Funktion des Systems
nachvollziehbar geprüft werden kann.

\subsection{Zielsetzung der
Testdurchführung}\label{zielsetzung-der-testdurchfuxfchrung}

Die Testdurchführung verfolgt folgende Ziele:

\begin{itemize}
\tightlist
\item
  Nachweis der Funktionsfähigkeit des Coasterbot-Prototyps
\item
  Überprüfung der autonomen Navigation
\item
  Validierung der Hinderniserkennung und Kollisionsvermeidung
\item
  Überprüfung der Aufnahme und Positionierung von Getränkeuntersetzern
\item
  Prüfung sicherheitskritischer Funktionen
\item
  Bewertung der Stabilität und Zuverlässigkeit der Software
\item
  Überprüfung der Simulationsumgebung hinsichtlich Reproduzierbarkeit
\end{itemize}

Die Tests dienen nicht nur der Fehlererkennung, sondern auch der
Bewertung der Systemqualität und der Erfüllung der Definition of Done.

\subsection{Geltungsbereich}\label{geltungsbereich}

Die Testspezifikation umfasst alle Funktionen, die im Rahmen des
Projekts umgesetzt werden.

Der Schwerpunkt liegt auf den Funktionen des Maturity Level 1:

\begin{itemize}
\tightlist
\item
  Navigation und Lokalisierung
\item
  Hinderniserkennung
\item
  Tischkantenerkennung
\item
  Sicherheitsfunktionen
\item
  Getränkeuntersetzeraufnahme
\item
  Getränkeuntersetzerplatzierung
\item
  Systemüberwachung
\item
  Simulation und Testbarkeit
\end{itemize}

Funktionen höherer Reifegrade werden nicht vollständig getestet, da sie
außerhalb des Projektumfangs liegen.

Dies betrifft insbesondere:

\begin{itemize}
\tightlist
\item
  automatisierten Getränketransport
\item
  Bestellaufnahme
\item
  Bezahlvorgänge
\item
  Reinigung
\item
  Mensch-Roboter-Interaktion
\item
  Netzwerkkommunikation
\end{itemize}

\subsection{Testobjekt}\label{testobjekt}

Das Testobjekt ist der Coasterbot-Prototyp, einschließlich seiner Hard-
und Softwarekomponenten.

Die Tests betrachten folgende Systembestandteile:

\subsubsection{Hardware}\label{hardware}

\begin{itemize}
\tightlist
\item
  Recheneinheit
\item
  Sensorik
\item
  Aktorik
\item
  Energieversorgung
\item
  Mechanische Konstruktion
\end{itemize}

\subsubsection{Software}\label{software}

\begin{itemize}
\tightlist
\item
  Hardwareabstraktion
\item
  Navigationskomponente
\item
  Lokalisierung
\item
  Bewegungssteuerung
\item
  Untersetzerhandling
\item
  Sicherheitslogik
\item
  Systemüberwachung
\end{itemize}

\subsubsection{Simulation}\label{simulation}

\begin{itemize}
\tightlist
\item
  Virtuelle Tischumgebung
\item
  Sensormodelle
\item
  Aktormodelle
\item
  Simulierte Fehlerzustände
\item
  Teststeuerung
\end{itemize}

\subsection{Testgrundlagen}\label{testgrundlagen}

Die Testfälle werden aus den Anforderungen des Lastenhefts sowie den
technischen Beschreibungen des Pflichtenhefts abgeleitet.

Jeder Testfall besitzt eine eindeutige Identifikation und beschreibt:

\begin{itemize}
\tightlist
\item
  Testziel
\item
  Bezug zu Anforderungen
\item
  Voraussetzungen
\item
  Testumgebung
\item
  Testablauf
\item
  Eingabedaten
\item
  Erwartetes Ergebnis
\item
  Tatsächliches Ergebnis
\item
  Bewertung
\end{itemize}

Die Tests werden so gestaltet, dass sie reproduzierbar durchgeführt
werden können.

\subsection{Teststrategie}\label{teststrategie}

Die Teststrategie folgt einem mehrstufigen Ansatz.

Die einzelnen Testebenen bauen aufeinander auf:

\subsection{Komponententests}\label{komponententests}

Komponententests prüfen einzelne Softwaremodule unabhängig voneinander.

Beispiele:

\begin{itemize}
\tightlist
\item
  Prüfung der Positionsberechnung
\item
  Prüfung von Zustandsübergängen
\item
  Prüfung der Sensorverarbeitung
\item
  Prüfung der Fehlererkennung
\end{itemize}

Ziel ist die frühzeitige Erkennung von Fehlern innerhalb einzelner
Komponenten.

\subsection{Integrationstests}\label{integrationstests}

Integrationstests prüfen die Kommunikation zwischen mehreren
Komponenten.

Beispiele:

\begin{itemize}
\tightlist
\item
  Übergabe von Sensordaten an die Navigation
\item
  Übergabe von Navigationsbefehlen an die Bewegungssteuerung
\item
  Kommunikation zwischen Sicherheitslogik und Aktorsteuerung
\end{itemize}

Ziel ist der Nachweis eines korrekten Zusammenspiels der
Systemkomponenten.

\subsection{Simulationstests}\label{simulationstests}

Simulationstests überprüfen das Verhalten des Gesamtsystems innerhalb
der virtuellen Umgebung.

Dabei werden insbesondere geprüft:

\begin{itemize}
\tightlist
\item
  Navigation
\item
  Hindernisvermeidung
\item
  Tischkantenerkennung
\item
  Untersetzerhandling
\item
  Fehlerzustände
\end{itemize}

Die Simulation ermöglicht reproduzierbare Tests ohne Abhängigkeit von
der physischen Hardware.

\subsection{Systemtests}\label{systemtests}

Systemtests überprüfen den vollständigen Coasterbot-Prototyp.

Dabei werden Hardware und Software gemeinsam betrachtet.

Beispiele:

\begin{itemize}
\tightlist
\item
  vollständiger Navigationsablauf
\item
  Aufnahme eines Untersetzers
\item
  Sicherheitsstopp
\item
  Verhalten bei Fehlerzuständen
\end{itemize}

\subsection{Testumgebung}\label{testumgebung}

Die Tests werden in mehreren Umgebungen durchgeführt.

\subsubsection{Entwicklungsumgebung}\label{entwicklungsumgebung}

Die Entwicklungsumgebung dient zur Durchführung automatisierter
Softwaretests.

Sie umfasst:

\begin{itemize}
\tightlist
\item
  Quellcodeverwaltung
\item
  Build-System
\item
  Testframework
\item
  Analysewerkzeuge
\end{itemize}

\subsubsection{Simulationsumgebung}\label{simulationsumgebung}

Die Simulationsumgebung ermöglicht die Prüfung des Systemverhaltens ohne
reale Hardware.

Sie beinhaltet:

\begin{itemize}
\tightlist
\item
  virtuelle Tischumgebung
\item
  simulierte Sensoren
\item
  simulierte Aktoren
\item
  definierte Testszenarien
\end{itemize}

\subsubsection{Physischer Prototyp}\label{physischer-prototyp}

Nach erfolgreicher Validierung in der Simulation werden ausgewählte
Tests auf dem realen Prototyp durchgeführt.

Dabei werden insbesondere geprüft:

\begin{itemize}
\tightlist
\item
  mechanische Funktion
\item
  reale Sensorik
\item
  reale Bewegung
\item
  Interaktion der Hardwarekomponenten
\end{itemize}

\subsection{Bewertung von
Testergebnissen}\label{bewertung-von-testergebnissen}

Ein Testfall wird anhand des Vergleichs zwischen erwartetem und
tatsächlichem Ergebnis bewertet.

Folgende Ergebnisse sind möglich:

\begin{longtable}[]{@{}ll@{}}
\toprule\noalign{}
Ergebnis & Bedeutung \\
\midrule\noalign{}
\endhead
\bottomrule\noalign{}
\endlastfoot
Bestanden & Erwartetes Ergebnis wurde erreicht \\
Fehlgeschlagen & Erwartetes Ergebnis wurde nicht erreicht \\
Nicht ausführbar & Voraussetzungen waren nicht erfüllt \\
\end{longtable}

Fehlgeschlagene Tests werden dokumentiert und führen zu einer
Fehleranalyse sowie gegebenenfalls zu einer Anpassung der
Implementierung.

\subsection{Aufbau der
Testspezifikation}\label{aufbau-der-testspezifikation}

Die weiteren Kapitel dieser Testspezifikation sind wie folgt aufgebaut:

\textbf{Kapitel 2 -- Testmethodik und Testmanagement} Beschreibung der
Testorganisation, Testprozesse und Dokumentation.

\textbf{Kapitel 3 -- Komponententests} Beschreibung der Tests einzelner
Softwaremodule.

\textbf{Kapitel 4 -- Integrationstests} Beschreibung der Tests zwischen
Softwarekomponenten.

\textbf{Kapitel 5 -- Simulationstests} Beschreibung der virtuellen
Testszenarien.

\textbf{Kapitel 6 -- Systemtests} Beschreibung der Tests am
vollständigen Prototyp.

\textbf{Kapitel 7 -- Testergebnisse und Bewertung} Dokumentation und
Auswertung der durchgeführten Tests.

\textbf{Kapitel 8 -- Zusammenfassung} Bewertung der Testergebnisse
hinsichtlich der Projektziele.

\subsection{Zusammenfassung}\label{zusammenfassung}

Dieses Kapitel definiert die Grundlagen für die systematische Prüfung
des Coasterbot-Prototyps.

Durch die Kombination aus Komponenten-, Integrations-, Simulations- und
Systemtests wird sichergestellt, dass die entwickelten Funktionen
nachvollziehbar bewertet werden können.

Die Testspezifikation bildet damit die Grundlage für den Nachweis, dass
der entwickelte Prototyp die Anforderungen des Lasten- und
Pflichtenhefts erfüllt.

\section{Testmethodik und
Testmanagement}\label{testmethodik-und-testmanagement}

\subsection{Ziel der Testmethodik}\label{ziel-der-testmethodik}

Die Testmethodik definiert die Vorgehensweise zur Planung, Durchführung
und Bewertung der Tests für den Coasterbot-Prototyp.

Ziel ist es, durch eine strukturierte und nachvollziehbare
Vorgehensweise sicherzustellen, dass relevante Anforderungen geprüft
werden. Die Testmethodik orientiert sich an den Prinzipien des
systematischen Softwaretests und berücksichtigt sowohl die
Besonderheiten eingebetteter Systeme als auch die Anforderungen an
autonome Robotersysteme.

Die Testmethodik verfolgt folgende Ziele:

\begin{itemize}
\tightlist
\item
  vollständige Abdeckung der relevanten Anforderungen
\item
  frühzeitige Erkennung von Fehlern
\item
  reproduzierbare Testergebnisse
\item
  nachvollziehbare Dokumentation
\item
  Verifikation der technischen Umsetzung
\end{itemize}

\subsection{Testvorgehensmodell}\label{testvorgehensmodell}

Die Entwicklung und Prüfung des Coasterbots erfolgt nach einem
inkrementellen Testvorgehen.

Dabei werden Funktionen schrittweise entwickelt und anschließend auf
verschiedenen Ebenen geprüft.

Der Testprozess besteht aus folgenden Phasen:

\begin{enumerate}
\def\labelenumi{\arabic{enumi}.}
\tightlist
\item
  Testplanung
\item
  Testentwurf
\item
  Testimplementierung
\item
  Testdurchführung
\item
  Fehleranalyse
\item
  Testabschluss
\end{enumerate}

\subsubsection{Testplanung}\label{testplanung}

In der Testplanung werden die Rahmenbedingungen für die Tests
festgelegt.

Dazu gehören:

\begin{itemize}
\tightlist
\item
  Definition der Testziele
\item
  Auswahl geeigneter Testverfahren
\item
  Festlegung der Testumgebung
\item
  Identifikation benötigter Testdaten
\item
  Zuordnung der Anforderungen zu Testfällen
\end{itemize}

Die Testplanung basiert auf der Traceability-Matrix aus dem
Pflichtenheft.

\subsubsection{Testentwurf}\label{testentwurf}

Im Testentwurf werden konkrete Testfälle aus den Anforderungen
abgeleitet.

Jeder Testfall erhält eine eindeutige Identifikation.

Die Testfälle werden nach folgendem Schema aufgebaut:

\begin{longtable}[]{@{}ll@{}}
\toprule\noalign{}
Feld & Beschreibung \\
\midrule\noalign{}
\endhead
\bottomrule\noalign{}
\endlastfoot
Test-ID & Eindeutige Kennung des Testfalls \\
Titel & Name des Tests \\
Ziel & Beschreibung der zu prüfenden Funktion \\
Anforderungen & Zugeordnete Anforderungen \\
Voraussetzungen & Notwendige Ausgangsbedingungen \\
Testdaten & Verwendete Eingaben und Parameter \\
Ablauf & Beschreibung der Testschritte \\
Erwartetes Ergebnis & Sollverhalten \\
Bewertung & Testergebnis \\
\end{longtable}

\subsection{Testebenen}\label{testebenen}

Die Tests werden in verschiedene Ebenen unterteilt. Jede Ebene besitzt
eine eigene Zielsetzung.

\subsubsection{Komponententest}\label{komponententest}

Komponententests überprüfen einzelne Softwaremodule unabhängig
voneinander.

Ziel:

\begin{itemize}
\tightlist
\item
  korrekte Funktion einzelner Komponenten
\item
  Prüfung interner Logik
\item
  Fehlerisolierung
\end{itemize}

Beispiele:

\begin{itemize}
\tightlist
\item
  Prüfung der Lokalisierungsberechnung
\item
  Prüfung eines Zustandsautomaten
\item
  Prüfung der Fehlererkennung
\end{itemize}

Komponententests werden bevorzugt automatisiert durchgeführt.

\subsubsection{Integrationstest}\label{integrationstest}

Integrationstests prüfen das Zusammenspiel mehrerer Komponenten.

Dabei werden insbesondere Schnittstellen überprüft.

Beispiele:

\begin{itemize}
\tightlist
\item
  Sensorabstraktion liefert korrekte Daten an Navigation
\item
  Navigation erzeugt gültige Steuerbefehle
\item
  Sicherheitskomponente kann Bewegungen stoppen
\end{itemize}

Ziel ist der Nachweis, dass die Komponenten gemeinsam korrekt
funktionieren.

\subsubsection{Simulationstest}\label{simulationstest}

Simulationstests prüfen das Verhalten des Gesamtsystems unter
kontrollierten Bedingungen.

Die Simulation ermöglicht:

\begin{itemize}
\tightlist
\item
  reproduzierbare Szenarien
\item
  sichere Prüfung kritischer Situationen
\item
  Untersuchung verschiedener Umgebungsbedingungen
\end{itemize}

Beispiele:

\begin{itemize}
\tightlist
\item
  Hindernis auf der Fahrstrecke
\item
  Verlust eines Sensorsignals
\item
  Annäherung an eine Tischkante
\end{itemize}

\subsubsection{Systemtest}\label{systemtest}

Systemtests prüfen den vollständigen Prototyp.

Dabei werden alle relevanten Komponenten gemeinsam betrachtet:

\begin{itemize}
\tightlist
\item
  Mechanik
\item
  Elektronik
\item
  Software
\item
  Sensorik
\item
  Aktorik
\end{itemize}

Ziel ist die Überprüfung des Gesamtsystems unter realitätsnahen
Bedingungen.

\subsection{Testarten}\label{testarten}

Neben den Testebenen werden unterschiedliche Testarten eingesetzt.

\subsubsection{Funktionstest}\label{funktionstest}

Funktionstests prüfen, ob das System die geforderten Funktionen korrekt
ausführt.

Beispiele:

\begin{itemize}
\tightlist
\item
  Navigation zu einem Zielpunkt
\item
  Aufnahme eines Getränkeuntersetzers
\item
  Erkennen eines Hindernisses
\end{itemize}

\subsubsection{Sicherheitstest}\label{sicherheitstest}

Sicherheitstests überprüfen sicherheitskritische Funktionen.

Geprüft werden:

\begin{itemize}
\tightlist
\item
  Not-Aus
\item
  Tischkantenerkennung
\item
  Kollisionsvermeidung
\item
  Fehlerreaktionen
\end{itemize}

Diese Tests besitzen eine hohe Priorität, da Fehler zu Schäden am System
oder in der Umgebung führen können.

\subsubsection{Belastungstest}\label{belastungstest}

Belastungstests untersuchen das Verhalten unter erhöhten Anforderungen.

Beispiele:

\begin{itemize}
\tightlist
\item
  längere Betriebsdauer
\item
  häufige Fahrbewegungen
\item
  wiederholte Aufnahmevorgänge
\end{itemize}

Ziel ist die Bewertung der Stabilität und Zuverlässigkeit.

\subsubsection{Regressionstest}\label{regressionstest}

Regressionstests stellen sicher, dass Änderungen an der Software keine
bereits funktionierenden Funktionen beeinträchtigen.

Nach jeder größeren Änderung werden relevante Tests erneut ausgeführt.

\subsection{Priorisierung der
Testfälle}\label{priorisierung-der-testfuxe4lle}

Nicht alle Anforderungen besitzen dieselbe Kritikalität. Daher werden
Testfälle priorisiert.

Die Priorisierung erfolgt anhand folgender Kriterien:

\begin{itemize}
\tightlist
\item
  Sicherheitsrelevanz
\item
  Bedeutung für die Kernfunktion
\item
  Abhängigkeit anderer Komponenten
\item
  Einfluss auf die Definition of Done
\end{itemize}

Die Prioritätsstufen sind:

\begin{longtable}[]{@{}ll@{}}
\toprule\noalign{}
Priorität & Bedeutung \\
\midrule\noalign{}
\endhead
\bottomrule\noalign{}
\endlastfoot
P1 & Kritische Funktion, muss zwingend erfüllt sein \\
P2 & Wesentliche Funktion für den Prototyp \\
P3 & Unterstützende oder optionale Funktion \\
\end{longtable}

\subsubsection{Priorität P1 -- Kritische
Tests}\label{priorituxe4t-p1-kritische-tests}

Folgende Funktionen besitzen höchste Priorität:

\begin{itemize}
\tightlist
\item
  Erkennung von Tischkanten
\item
  Sicherheitsstopp
\item
  Kollisionsvermeidung
\item
  Grundlegende Bewegungsfähigkeit
\item
  Fehlererkennung
\end{itemize}

Ein Fehler in diesen Bereichen verhindert einen sicheren Betrieb des
Roboters.

\subsubsection{Priorität P2 --
Kernfunktionen}\label{priorituxe4t-p2-kernfunktionen}

Folgende Funktionen sind für die eigentliche Aufgabe des Prototyps
erforderlich:

\begin{itemize}
\tightlist
\item
  Navigation (zu Zielpunkten)
\item
  Lokalisierung
\item
  Untersetzeraufnahme
\item
  Untersetzerplatzierung
\item
  Systemüberwachung
\end{itemize}

\subsubsection{Priorität P3 --
Erweiterungen}\label{priorituxe4t-p3-erweiterungen}

Diese Funktionen werden nachrangig betrachtet:

\begin{itemize}
\tightlist
\item
  (Erweiterte) Simulationen
\item
  Optimierungen
\item
  Vorbereitung zukünftiger Ausbaustufen
\end{itemize}

\subsection{Testdaten und
Testszenarien}\label{testdaten-und-testszenarien}

Die Testdaten werden so gewählt, dass typische sowie kritische
Situationen abgedeckt werden.

Beispiele:

\subsubsection{Normalszenario}\label{normalszenario}

\begin{itemize}
\tightlist
\item
  Freie Tischfläche
\item
  Keine Hindernisse
\item
  Ausreichender Energiezustand
\end{itemize}

Ziel:

Überprüfung des normalen Betriebs.

\subsubsection{Hindernisszenario}\label{hindernisszenario}

\begin{itemize}
\tightlist
\item
  Objekt befindet sich auf geplanter Route
\end{itemize}

Ziel:

Prüfung der Hinderniserkennung und Navigation.

\subsubsection{Grenzszenario}\label{grenzszenario}

\begin{itemize}
\tightlist
\item
  Roboter nähert sich Tischkante
\item
  Sensor liefert kritische Werte
\end{itemize}

Ziel:

Überprüfung der Sicherheitsfunktionen.

\subsubsection{Fehlerszenario}\label{fehlerszenario}

\begin{itemize}
\tightlist
\item
  Sensor fällt aus
\item
  Aktor reagiert nicht
\end{itemize}

Ziel:

Überprüfung der Fehlerbehandlung.

\subsection{Testautomatisierung}\label{testautomatisierung}

Ein wesentlicher Bestandteil des Testkonzepts ist die Automatisierung
wiederkehrender Prüfungen.

Automatisierte Tests ermöglichen:

\begin{itemize}
\tightlist
\item
  schnelle Wiederholung von Testfällen
\item
  reproduzierbare Ergebnisse
\item
  frühzeitige Fehlererkennung
\item
  Unterstützung der Softwareentwicklung
\end{itemize}

Automatisiert werden insbesondere:

\begin{itemize}
\tightlist
\item
  Komponententests
\item
  Schnittstellentests
\item
  Simulationstests
\item
  Zustandsübergänge
\end{itemize}

\subsection{Fehlerverwaltung}\label{fehlerverwaltung}

Festgestellte Fehler werden strukturiert dokumentiert.

Jeder Fehler enthält:

\begin{longtable}[]{@{}ll@{}}
\toprule\noalign{}
Feld & Beschreibung \\
\midrule\noalign{}
\endhead
\bottomrule\noalign{}
\endlastfoot
Fehler-ID & Eindeutige Kennung \\
Beschreibung & Fehlerbeschreibung \\
Betroffene Komponente & Verantwortliches Modul \\
Priorität & Kritikalität \\
Reproduzierbarkeit & Möglichkeit der Wiederholung \\
Status & Bearbeitungsstand \\
\end{longtable}

Mögliche Fehlerstatus:

\begin{longtable}[]{@{}ll@{}}
\toprule\noalign{}
Status & Bedeutung \\
\midrule\noalign{}
\endhead
\bottomrule\noalign{}
\endlastfoot
Neu & Fehler wurde erkannt \\
Analyse & Ursache wird untersucht \\
Behoben & Fehler wurde korrigiert \\
Verifiziert & Korrektur wurde getestet \\
\end{longtable}

\subsection{Testabschluss}\label{testabschluss}

Ein Testzyklus gilt als abgeschlossen, wenn:

\begin{itemize}
\tightlist
\item
  alle geplanten Tests durchgeführt wurden
\item
  Ergebnisse dokumentiert wurden
\item
  kritische Fehler behoben wurden
\item
  die Testabdeckung bewertet wurde
\end{itemize}

Die Ergebnisse werden in einem Testbericht zusammengefasst.

Dieser enthält:

\begin{itemize}
\tightlist
\item
  Anzahl ausgeführter Tests
\item
  erfolgreiche Tests
\item
  fehlgeschlagene Tests
\item
  offene Probleme
\item
  Bewertung des Systemzustands
\end{itemize}

\subsection{Zusammenfassung}\label{zusammenfassung-1}

Die Testmethodik definiert einen strukturierten Prozess zur Überprüfung
des Coasterbot-Prototyps.

Durch die Kombination verschiedener Testebenen und Testarten wird
sichergestellt, dass sowohl einzelne Softwarekomponenten als auch das
Gesamtsystem überprüft werden.

Die Priorisierung der Tests stellt sicher, dass sicherheitskritische und
für die Kernfunktion notwendige Anforderungen zuerst validiert werden.

Damit bildet dieses Testmanagement die Grundlage für die nachfolgenden
detaillierten Testfälle.

\section{Komponententests}\label{komponententests-1}

\subsection{Ziel der Komponententests}\label{ziel-der-komponententests}

Komponententests überprüfen einzelne Softwaremodule des
Coasterbot-Systems unabhängig von anderen Systembestandteilen. Ziel ist
es, die korrekte Funktion der einzelnen Komponenten sicherzustellen,
bevor diese in das Gesamtsystem integriert werden.

Durch die isolierte Prüfung einzelner Komponenten können Fehler
frühzeitig erkannt und deren Ursachen eindeutig zugeordnet werden.

Die Komponententests konzentrieren sich auf die im Pflichtenheft
definierten Softwaremodule:

\begin{itemize}
\tightlist
\item
  Hardwareabstraktion
\item
  Lokalisierung
\item
  Navigation
\item
  Bewegungssteuerung
\item
  Untersetzerhandling
\item
  Sicherheitskomponente
\item
  Systemüberwachung
\end{itemize}

Die Tests werden bevorzugt automatisiert ausgeführt, um eine
wiederholbare und effiziente Qualitätssicherung zu ermöglichen.

\subsection{Testumgebung für
Komponententests}\label{testumgebung-fuxfcr-komponententests}

Die Komponententests werden ohne vollständige Hardware durchgeführt.
Hardwareabhängige Funktionen werden durch simulierte Schnittstellen
ersetzt.

Die Testumgebung besteht aus:

\begin{itemize}
\tightlist
\item
  Softwarekomponenten des Coasterbots
\item
  Simulierten Sensordaten
\item
  Simulierten Aktorzuständen
\item
  Testframework
\item
  Testdatenverwaltung
\end{itemize}

Die Hardwareabstraktion ermöglicht es, einzelne Komponenten unabhängig
von realen Sensoren und Aktoren zu prüfen.

\subsection{Testfallstruktur}\label{testfallstruktur}

Alle Komponententests werden nach einer einheitlichen Struktur
dokumentiert.

Jeder Testfall enthält:

\begin{longtable}[]{@{}ll@{}}
\toprule\noalign{}
Feld & Beschreibung \\
\midrule\noalign{}
\endhead
\bottomrule\noalign{}
\endlastfoot
Test-ID & Eindeutige Kennung \\
Komponente & Getestetes Modul \\
Ziel & Zweck des Tests \\
Anforderungen & Zugeordnete Lastenheft-Anforderungen \\
Voraussetzungen & Notwendige Bedingungen \\
Eingaben & Testdaten \\
Ablauf & Durchführung \\
Erwartetes Ergebnis & Sollverhalten \\
Bewertung & Testergebnis \\
\end{longtable}

\subsection{Komponententest:
Hardwareabstraktion}\label{komponententest-hardwareabstraktion}

\subsubsection{Testfall CT-HW-001: Verarbeitung von
Sensordaten}\label{testfall-ct-hw-001-verarbeitung-von-sensordaten}

\paragraph{Ziel}\label{ziel}

Überprüfung, ob die Hardwareabstraktion Sensordaten korrekt einliest und
für die darüberliegenden Softwarekomponenten bereitstellt.

\paragraph{Zugeordnete Anforderungen}\label{zugeordnete-anforderungen}

\begin{itemize}
\tightlist
\item
  NFR-008: Software muss Hardwareabstraktion unterstützen
\item
  NFR-007: Sensoren und Aktoren müssen austauschbar sein
\end{itemize}

\paragraph{Voraussetzungen}\label{voraussetzungen}

\begin{itemize}
\tightlist
\item
  Simulierter Sensor ist verfügbar
\item
  Hardwareabstraktionsschnittstelle ist initialisiert
\end{itemize}

\paragraph{Eingaben}\label{eingaben}

Simulierter Abstandswert:

\begin{verbatim}
Entfernung = 500 mm
\end{verbatim}

\paragraph{Ablauf}\label{ablauf}

\begin{enumerate}
\def\labelenumi{\arabic{enumi}.}
\tightlist
\item
  Simulierter Sensor liefert Messwert
\item
  Hardwareabstraktion empfängt den Messwert
\item
  Daten werden in internes Format übertragen
\item
  Navigationskomponente ruft Sensordaten ab
\end{enumerate}

\paragraph{Erwartetes Ergebnis}\label{erwartetes-ergebnis}

\begin{itemize}
\tightlist
\item
  Sensordaten werden korrekt verarbeitet
\item
  Datenformat entspricht der definierten Schnittstelle
\item
  Navigation erhält gültige Informationen
\end{itemize}

\paragraph{Bewertung}\label{bewertung}

Bestanden, wenn die übertragenen Daten mit den Eingangsdaten
übereinstimmen.

\subsection{Komponententest:
Lokalisierung}\label{komponententest-lokalisierung}

\subsubsection{Testfall CT-LOC-001:
Positionsaktualisierung}\label{testfall-ct-loc-001-positionsaktualisierung}

\paragraph{Ziel}\label{ziel-1}

Überprüfung der Berechnung und Aktualisierung der Roboterposition.

\paragraph{Zugeordnete Anforderungen}\label{zugeordnete-anforderungen-1}

\begin{itemize}
\tightlist
\item
  NAV-006: Der Bot muss seine Position bestimmen können
\item
  NAV-007: Der Bot muss seine Orientierung bestimmen können
\end{itemize}

\paragraph{Voraussetzungen}\label{voraussetzungen-1}

\begin{itemize}
\tightlist
\item
  Lokalisierungskomponente ist aktiv
\item
  Startposition ist definiert
\end{itemize}

\paragraph{Eingaben}\label{eingaben-1}

Startposition:

\[
(x,y,\theta)=(0,0,0)
\]

Bewegungsänderung:

\[
\Delta x=100mm
\]

\paragraph{Ablauf}\label{ablauf-1}

\begin{enumerate}
\def\labelenumi{\arabic{enumi}.}
\tightlist
\item
  Startposition wird gesetzt
\item
  Bewegungsinformation wird übergeben
\item
  Lokalisierung verarbeitet die Daten
\item
  Neue Position wird berechnet
\end{enumerate}

\paragraph{Erwartetes Ergebnis}\label{erwartetes-ergebnis-1}

Die Position wird korrekt aktualisiert:

\[
(x,y,\theta)=(100,0,0)
\]

\paragraph{Bewertung}\label{bewertung-1}

Bestanden, wenn die berechnete Position innerhalb der definierten
Genauigkeit liegt.

\subsection{Komponententest:
Navigation}\label{komponententest-navigation}

\subsubsection{Testfall CT-NAV-001: Berechnung einer gültigen
Route}\label{testfall-ct-nav-001-berechnung-einer-guxfcltigen-route}

\paragraph{Ziel}\label{ziel-2}

Überprüfung, ob die Navigationskomponente eine gültige Route zwischen
Start- und Zielpunkt erzeugen kann.

\paragraph{Zugeordnete Anforderungen}\label{zugeordnete-anforderungen-2}

\begin{itemize}
\tightlist
\item
  NAV-005: Der Bot muss seine Fahrtroute dynamisch anpassen
\item
  NAV-008: Der Bot muss seinen Zielpunkt erreichen
\end{itemize}

\paragraph{Voraussetzungen}\label{voraussetzungen-2}

\begin{itemize}
\tightlist
\item
  Startposition bekannt
\item
  Zielposition bekannt
\item
  Keine Hindernisse vorhanden
\end{itemize}

\paragraph{Eingaben}\label{eingaben-2}

Start:

\[
(0,0)
\]

Ziel:

\[
(1000,500)
\]

\paragraph{Ablauf}\label{ablauf-2}

\begin{enumerate}
\def\labelenumi{\arabic{enumi}.}
\tightlist
\item
  Startposition wird gesetzt
\item
  Zielposition wird übergeben
\item
  Navigation berechnet Route
\item
  Route wird ausgegeben
\end{enumerate}

\paragraph{Erwartetes Ergebnis}\label{erwartetes-ergebnis-2}

\begin{itemize}
\tightlist
\item
  Eine gültige Route wird erzeugt
\item
  Zielpunkt ist erreichbar
\item
  Route enthält keine ungültigen Bewegungen
\end{itemize}

\paragraph{Bewertung}\label{bewertung-2}

Bestanden, wenn eine gültige Bewegungssequenz erzeugt wird.

\subsubsection{Testfall CT-NAV-002:
Hinderniserkennung}\label{testfall-ct-nav-002-hinderniserkennung}

\paragraph{Ziel}\label{ziel-3}

Überprüfung der Reaktion der Navigation auf ein erkanntes Hindernis.

\paragraph{Zugeordnete Anforderungen}\label{zugeordnete-anforderungen-3}

\begin{itemize}
\tightlist
\item
  NAV-001: Hindernisse erkennen
\item
  NAV-002: Hindernissen ausweichen
\end{itemize}

\paragraph{Voraussetzungen}\label{voraussetzungen-3}

\begin{itemize}
\tightlist
\item
  Navigationsziel ist gesetzt
\item
  Hindernisdaten werden simuliert
\end{itemize}

\paragraph{Eingaben}\label{eingaben-3}

Hindernisposition:

\[
(500,0)
\]

\paragraph{Ablauf}\label{ablauf-3}

\begin{enumerate}
\def\labelenumi{\arabic{enumi}.}
\tightlist
\item
  Navigation startet
\item
  Hindernis wird erkannt
\item
  Route wird neu berechnet
\end{enumerate}

\paragraph{Erwartetes Ergebnis}\label{erwartetes-ergebnis-3}

\begin{itemize}
\tightlist
\item
  Hindernis wird erkannt
\item
  Neue Route vermeidet das Hindernis
\item
  Keine Kollision entsteht
\end{itemize}

\paragraph{Bewertung}\label{bewertung-3}

Bestanden, wenn eine kollisionsfreie Route erzeugt wird.

\subsection{Komponententest:
Bewegungssteuerung}\label{komponententest-bewegungssteuerung}

\subsubsection{Testfall CT-MOT-001: Umsetzung von
Bewegungsbefehlen}\label{testfall-ct-mot-001-umsetzung-von-bewegungsbefehlen}

\paragraph{Ziel}\label{ziel-4}

Überprüfung, ob Bewegungsbefehle korrekt in Aktorbefehle umgesetzt
werden.

\paragraph{Zugeordnete Anforderungen}\label{zugeordnete-anforderungen-4}

\begin{itemize}
\tightlist
\item
  NAV-009: Der Bot muss sich kontrolliert bewegen
\item
  SAF-004: Der Bot darf nicht mit hoher Geschwindigkeit kollidieren
\end{itemize}

\paragraph{Voraussetzungen}\label{voraussetzungen-4}

\begin{itemize}
\tightlist
\item
  Motorsteuerung ist initialisiert
\end{itemize}

\paragraph{Eingaben}\label{eingaben-4}

Bewegungsbefehl:

\begin{verbatim}
Geschwindigkeit: 100 mm/s
Richtung: vorwärts
\end{verbatim}

\paragraph{Ablauf}\label{ablauf-4}

\begin{enumerate}
\def\labelenumi{\arabic{enumi}.}
\tightlist
\item
  Bewegungsbefehl wird übergeben
\item
  Steuerung berechnet Aktorsignale
\item
  Motorzustand wird geprüft
\end{enumerate}

\paragraph{Erwartetes Ergebnis}\label{erwartetes-ergebnis-4}

\begin{itemize}
\tightlist
\item
  Motoren erhalten korrekte Steuerwerte
\item
  Geschwindigkeit entspricht Vorgabe
\item
  Sicherheitsgrenzen werden eingehalten
\end{itemize}

\paragraph{Bewertung}\label{bewertung-4}

Bestanden, wenn die Ausgangswerte korrekt erzeugt werden.

\subsection{Komponententest:
Untersetzerhandling}\label{komponententest-untersetzerhandling}

\subsubsection{Testfall CT-CST-001: Aufnahme eines
Getränkeuntersetzers}\label{testfall-ct-cst-001-aufnahme-eines-getruxe4nkeuntersetzers}

\paragraph{Ziel}\label{ziel-5}

Überprüfung des Aufnahmeprozesses.

\paragraph{Zugeordnete Anforderungen}\label{zugeordnete-anforderungen-5}

\begin{itemize}
\tightlist
\item
  CST-001: Untersetzer aufnehmen
\item
  CST-005: Erfolgreiche Aufnahme erkennen
\end{itemize}

\paragraph{Voraussetzungen}\label{voraussetzungen-5}

\begin{itemize}
\tightlist
\item
  Untersetzer befindet sich an definierter Position
\item
  Aufnahmemechanismus ist verfügbar
\end{itemize}

\paragraph{Eingaben}\label{eingaben-5}

Untersetzerstatus:

\begin{verbatim}
Verfügbar = TRUE
\end{verbatim}

\paragraph{Ablauf}\label{ablauf-5}

\begin{enumerate}
\def\labelenumi{\arabic{enumi}.}
\tightlist
\item
  Aufnahmeprozess wird gestartet
\item
  Aktor wird aktiviert
\item
  Sensorstatus wird geprüft
\end{enumerate}

\paragraph{Erwartetes Ergebnis}\label{erwartetes-ergebnis-5}

\begin{itemize}
\tightlist
\item
  Aufnahme wird durchgeführt
\item
  Erfolgreiche Aufnahme wird erkannt
\item
  Zustand wechselt zu ``Untersetzer aufgenommen''
\end{itemize}

\paragraph{Bewertung}\label{bewertung-5}

Bestanden, wenn der Zustand korrekt aktualisiert wird.

\subsubsection{Testfall CT-CST-002: Fehler bei
Untersetzeraufnahme}\label{testfall-ct-cst-002-fehler-bei-untersetzeraufnahme}

\paragraph{Ziel}\label{ziel-6}

Überprüfung der Fehlererkennung während der Aufnahme.

\paragraph{Zugeordnete Anforderungen}\label{zugeordnete-anforderungen-6}

\begin{itemize}
\tightlist
\item
  CST-005: Aufnahme erkennen
\item
  MON-005: Fehlgeschlagene Aufgaben melden
\end{itemize}

\paragraph{Voraussetzungen}\label{voraussetzungen-6}

\begin{itemize}
\tightlist
\item
  Aufnahmevorgang gestartet
\end{itemize}

\paragraph{Eingaben}\label{eingaben-6}

Simulierter Fehler:

\begin{verbatim}
Untersetzer erkannt = FALSE
\end{verbatim}

\paragraph{Ablauf}\label{ablauf-6}

\begin{enumerate}
\def\labelenumi{\arabic{enumi}.}
\tightlist
\item
  Aufnahme wird gestartet
\item
  Sensor meldet fehlende Aufnahme
\item
  System verarbeitet Fehler
\end{enumerate}

\paragraph{Erwartetes Ergebnis}\label{erwartetes-ergebnis-6}

\begin{itemize}
\tightlist
\item
  Fehler wird erkannt
\item
  Aufgabe wird abgebrochen
\item
  Fehler wird gemeldet
\end{itemize}

\paragraph{Bewertung}\label{bewertung-6}

Bestanden, wenn der Fehler korrekt verarbeitet wird.

\subsection{Komponententest:
Sicherheitskomponente}\label{komponententest-sicherheitskomponente}

\subsubsection{Testfall CT-SAF-001:
Sicherheitsstopp}\label{testfall-ct-saf-001-sicherheitsstopp}

\paragraph{Ziel}\label{ziel-7}

Überprüfung der Reaktion auf eine Gefahrensituation.

\paragraph{Zugeordnete Anforderungen}\label{zugeordnete-anforderungen-7}

\begin{itemize}
\tightlist
\item
  SAF-001: Bei Gefahr sofort anhalten
\item
  SAF-006: Not-Aus unterstützen
\end{itemize}

\paragraph{Voraussetzungen}\label{voraussetzungen-7}

\begin{itemize}
\tightlist
\item
  Roboter befindet sich im Bewegungszustand
\end{itemize}

\paragraph{Eingaben}\label{eingaben-7}

Not-Aus-Signal:

\begin{verbatim}
EmergencyStop = TRUE
\end{verbatim}

\paragraph{Ablauf}\label{ablauf-7}

\begin{enumerate}
\def\labelenumi{\arabic{enumi}.}
\tightlist
\item
  Bewegungssteuerung ist aktiv
\item
  Sicherheitskomponente erhält Stoppsignal
\item
  Sicherheitszustand wird aktiviert
\end{enumerate}

\paragraph{Erwartetes Ergebnis}\label{erwartetes-ergebnis-7}

\begin{itemize}
\tightlist
\item
  Bewegungsbefehle werden unterbrochen
\item
  Roboter wechselt in sicheren Zustand
\item
  Status wird aktualisiert
\end{itemize}

\paragraph{Bewertung}\label{bewertung-7}

Bestanden, wenn der Stopp innerhalb der definierten Reaktionszeit
erfolgt.

\subsection{Komponententest:
Systemüberwachung}\label{komponententest-systemuxfcberwachung}

\subsubsection{Testfall CT-MON-001: Fehlererkennung und
Protokollierung}\label{testfall-ct-mon-001-fehlererkennung-und-protokollierung}

\paragraph{Ziel}\label{ziel-8}

Überprüfung der Diagnosefunktion.

\paragraph{Zugeordnete Anforderungen}\label{zugeordnete-anforderungen-8}

\begin{itemize}
\tightlist
\item
  MON-001: Fehler erkennen
\item
  MON-002: Fehler protokollieren
\end{itemize}

\paragraph{Voraussetzungen}\label{voraussetzungen-8}

\begin{itemize}
\tightlist
\item
  Systemüberwachung ist aktiv
\end{itemize}

\paragraph{Eingaben}\label{eingaben-8}

Simulierter Fehler:

\begin{verbatim}
Sensorfehler = TRUE
\end{verbatim}

\paragraph{Ablauf}\label{ablauf-8}

\begin{enumerate}
\def\labelenumi{\arabic{enumi}.}
\tightlist
\item
  Fehler wird erzeugt
\item
  Systemüberwachung verarbeitet Ereignis
\item
  Fehler wird gespeichert
\end{enumerate}

\paragraph{Erwartetes Ergebnis}\label{erwartetes-ergebnis-8}

\begin{itemize}
\tightlist
\item
  Fehler wird erkannt
\item
  Fehler wird protokolliert
\item
  Status wird aktualisiert
\end{itemize}

\paragraph{Bewertung}\label{bewertung-8}

Bestanden, wenn der Fehler vollständig dokumentiert wurde.

\subsection{Zusammenfassung der
Komponententests}\label{zusammenfassung-der-komponententests}

Die Komponententests überprüfen die grundlegenden Funktionen der
einzelnen Softwaremodule des Coasterbots.

Durch die isolierte Prüfung werden Fehler frühzeitig erkannt und die
Grundlage für nachfolgende Integrations- und Systemtests geschaffen.

Die wichtigsten überprüften Bereiche sind:

\begin{itemize}
\tightlist
\item
  Hardwareabstraktion
\item
  Navigation
\item
  Lokalisierung
\item
  Bewegungssteuerung
\item
  Untersetzerhandling
\item
  Sicherheit
\item
  Systemüberwachung
\end{itemize}

Die erfolgreichen Komponententests bilden die Voraussetzung für die
Integration der einzelnen Softwarebestandteile zu einem vollständigen
Robotersystem.

\section{Integrationstests}\label{integrationstests-1}

\subsection{Ziel der
Integrationstests}\label{ziel-der-integrationstests}

Integrationstests überprüfen das Zusammenspiel mehrerer
Softwarekomponenten des Coasterbot-Systems. Während Komponententests
einzelne Module isoliert betrachten, wird bei Integrationstests die
korrekte Kommunikation und Zusammenarbeit zwischen den Komponenten
untersucht.

Ziel der Integrationstests ist es, sicherzustellen, dass die einzelnen
Systembestandteile gemeinsam die geforderten Funktionen erfüllen.

Der Fokus liegt dabei auf:

\begin{itemize}
\tightlist
\item
  korrekter Datenübertragung zwischen Komponenten
\item
  definierten Schnittstellen
\item
  korrekten Zustandsübergängen
\item
  fehlerfreier Zusammenarbeit von Hard- und Softwareabstraktion
\item
  sicherem Verhalten bei kritischen Situationen
\end{itemize}

Die Integrationstests bilden die Grundlage für die anschließenden
Systemtests, bei denen der vollständige Prototyp betrachtet wird.

\subsection{Testumgebung für
Integrationstests}\label{testumgebung-fuxfcr-integrationstests}

Die Integrationstests werden in einer Umgebung durchgeführt, die mehrere
reale oder simulierte Komponenten kombiniert.

Die Testumgebung umfasst:

\begin{itemize}
\tightlist
\item
  Softwarekomponenten des Coasterbots
\item
  Hardwareabstraktionsschicht
\item
  simulierte Sensordaten
\item
  simulierte Aktorzustände
\item
  Simulationsumgebung
\item
  Testframework
\end{itemize}

Je nach Testfall können einzelne Hardwarekomponenten durch
Simulationsmodelle ersetzt werden.

Dadurch kann das Verhalten der Software unabhängig von der finalen
Hardware überprüft werden.

\subsection{Testfallstruktur}\label{testfallstruktur-1}

Die Integrationstestfälle werden nach folgendem Schema beschrieben:

\begin{longtable}[]{@{}ll@{}}
\toprule\noalign{}
Feld & Beschreibung \\
\midrule\noalign{}
\endhead
\bottomrule\noalign{}
\endlastfoot
Test-ID & Eindeutige Identifikation \\
Testname & Bezeichnung des Tests \\
Ziel & Zu überprüfende Integration \\
Anforderungen & Zugeordnete Anforderungen \\
Beteiligte Komponenten & Verwendete Systemmodule \\
Voraussetzungen & Ausgangszustand \\
Ablauf & Testschritte \\
Erwartetes Ergebnis & Sollverhalten \\
Bewertung & Testergebnis \\
\end{longtable}

\subsection{Integrationstest: Sensorik und
Navigation}\label{integrationstest-sensorik-und-navigation}

\subsubsection{Testfall IT-SEN-NAV-001: Verarbeitung von
Hindernisdaten}\label{testfall-it-sen-nav-001-verarbeitung-von-hindernisdaten}

\paragraph{Ziel}\label{ziel-9}

Überprüfung der korrekten Übergabe von Sensordaten an die
Navigationskomponente.

\paragraph{Zugeordnete Anforderungen}\label{zugeordnete-anforderungen-9}

\begin{itemize}
\tightlist
\item
  NAV-001: Der Bot muss Hindernisse erkennen
\item
  NAV-002: Der Bot muss Hindernissen ausweichen
\item
  NFR-008: Software muss Hardwareabstraktion unterstützen
\end{itemize}

\paragraph{Beteiligte Komponenten}\label{beteiligte-komponenten}

\begin{itemize}
\tightlist
\item
  Sensorabstraktion
\item
  Hinderniserkennung
\item
  Navigation
\end{itemize}

\paragraph{Voraussetzungen}\label{voraussetzungen-9}

\begin{itemize}
\tightlist
\item
  Navigation ist aktiv
\item
  Simulierter Abstandssensor ist verfügbar
\item
  Eine Route zu einem Zielpunkt wurde erstellt
\end{itemize}

\paragraph{Eingaben}\label{eingaben-9}

Sensordaten:

\begin{verbatim}
Hindernis erkannt = TRUE
Abstand = 200 mm
Position = (500,300)
\end{verbatim}

\paragraph{Ablauf}\label{ablauf-9}

\begin{enumerate}
\def\labelenumi{\arabic{enumi}.}
\tightlist
\item
  Sensorabstraktion empfängt Messwerte
\item
  Messwerte werden an die Hinderniserkennung übertragen
\item
  Hinderniserkennung bewertet die Situation
\item
  Navigationskomponente erhält die Information über das Hindernis
\item
  Neue Route wird berechnet
\end{enumerate}

\paragraph{Erwartetes Ergebnis}\label{erwartetes-ergebnis-9}

\begin{itemize}
\tightlist
\item
  Sensordaten werden korrekt übertragen
\item
  Hindernis wird erkannt
\item
  Navigation reagiert auf die Änderung
\item
  Neue kollisionsfreie Route wird erzeugt
\end{itemize}

\paragraph{Bewertung}\label{bewertung-9}

Bestanden, wenn die Navigation innerhalb der definierten Reaktionszeit
eine gültige Anpassung erzeugt.

\subsection{Integrationstest: Lokalisierung und
Navigation}\label{integrationstest-lokalisierung-und-navigation}

\subsubsection{Testfall IT-LOC-NAV-001: Zielnavigation mit
Positionsdaten}\label{testfall-it-loc-nav-001-zielnavigation-mit-positionsdaten}

\paragraph{Ziel}\label{ziel-10}

Überprüfung, ob die Navigation die aktuelle Position korrekt verwendet.

\paragraph{Zugeordnete
Anforderungen}\label{zugeordnete-anforderungen-10}

\begin{itemize}
\tightlist
\item
  NAV-006: Der Bot muss seine Position bestimmen können
\item
  NAV-007: Der Bot muss seine Orientierung bestimmen können
\item
  NAV-008: Der Bot muss seinen Zielpunkt erreichen
\end{itemize}

\paragraph{Beteiligte Komponenten}\label{beteiligte-komponenten-1}

\begin{itemize}
\tightlist
\item
  Lokalisierung
\item
  Navigation
\item
  Bewegungssteuerung
\end{itemize}

\paragraph{Voraussetzungen}\label{voraussetzungen-10}

\begin{itemize}
\tightlist
\item
  Roboterposition ist initialisiert
\item
  Zielpunkt ist bekannt
\end{itemize}

\paragraph{Eingaben}\label{eingaben-10}

Startposition:

\begin{verbatim}
x = 0 mm
y = 0 mm
Orientierung = 0°
\end{verbatim}

Ziel:

\begin{verbatim}
x = 1000 mm
y = 500 mm
\end{verbatim}

\paragraph{Ablauf}\label{ablauf-10}

\begin{enumerate}
\def\labelenumi{\arabic{enumi}.}
\tightlist
\item
  Lokalisierung stellt aktuelle Position bereit
\item
  Navigation berechnet Bewegungsrichtung
\item
  Bewegungssteuerung erhält Fahrbefehle
\item
  Positionsänderungen werden zurückgemeldet
\item
  Navigation aktualisiert die Route
\end{enumerate}

\paragraph{Erwartetes Ergebnis}\label{erwartetes-ergebnis-10}

\begin{itemize}
\tightlist
\item
  Navigation verwendet aktuelle Positionsdaten
\item
  Fahrbefehle werden korrekt erzeugt
\item
  Roboter bewegt sich in Richtung Ziel
\item
  Ziel wird innerhalb der definierten Genauigkeit erreicht
\end{itemize}

\paragraph{Bewertung}\label{bewertung-10}

Bestanden, wenn die Positionsabweichung kleiner als der definierte
Grenzwert ist.

\subsection{Integrationstest: Navigation und
Bewegungssteuerung}\label{integrationstest-navigation-und-bewegungssteuerung}

\subsubsection{Testfall IT-NAV-MOT-001: Umsetzung von
Navigationsbefehlen}\label{testfall-it-nav-mot-001-umsetzung-von-navigationsbefehlen}

\paragraph{Ziel}\label{ziel-11}

Überprüfung der Schnittstelle zwischen Navigation und Motorsteuerung.

\paragraph{Zugeordnete
Anforderungen}\label{zugeordnete-anforderungen-11}

\begin{itemize}
\tightlist
\item
  NAV-009: Der Bot muss sich kontrolliert bewegen
\item
  SAF-004: Der Bot darf nicht mit hoher Geschwindigkeit in Hindernisse
  fahren
\end{itemize}

\paragraph{Beteiligte Komponenten}\label{beteiligte-komponenten-2}

\begin{itemize}
\tightlist
\item
  Navigation
\item
  Bewegungssteuerung
\item
  Hardwareabstraktion
\end{itemize}

\paragraph{Voraussetzungen}\label{voraussetzungen-11}

\begin{itemize}
\tightlist
\item
  Route wurde berechnet
\item
  Motorsteuerung ist betriebsbereit
\end{itemize}

\paragraph{Eingaben}\label{eingaben-11}

Navigationsbefehl:

\begin{verbatim}
Geschwindigkeit = 100 mm/s
Richtung = 45°
Distanz = 500 mm
\end{verbatim}

\paragraph{Ablauf}\label{ablauf-11}

\begin{enumerate}
\def\labelenumi{\arabic{enumi}.}
\tightlist
\item
  Navigation erzeugt Bewegungsbefehl
\item
  Bewegungssteuerung verarbeitet den Befehl
\item
  Aktorschnittstelle erzeugt Steuersignale
\item
  Bewegungsstatus wird zurückgemeldet
\end{enumerate}

\paragraph{Erwartetes Ergebnis}\label{erwartetes-ergebnis-11}

\begin{itemize}
\tightlist
\item
  Bewegungsbefehle werden korrekt umgesetzt
\item
  Geschwindigkeit bleibt innerhalb definierter Grenzen
\item
  Bewegungsstatus wird aktualisiert
\end{itemize}

\paragraph{Bewertung}\label{bewertung-11}

Bestanden, wenn die erzeugte Bewegung dem Navigationsauftrag entspricht.

\subsection{Integrationstest: Sicherheitslogik und
Bewegungssteuerung}\label{integrationstest-sicherheitslogik-und-bewegungssteuerung}

\subsubsection{Testfall IT-SAF-MOT-001: Unterbrechen einer Bewegung bei
Gefahr}\label{testfall-it-saf-mot-001-unterbrechen-einer-bewegung-bei-gefahr}

\paragraph{Ziel}\label{ziel-12}

Überprüfung, ob die Sicherheitskomponente jederzeit Einfluss auf die
Bewegungssteuerung nehmen kann.

\paragraph{Zugeordnete
Anforderungen}\label{zugeordnete-anforderungen-12}

\begin{itemize}
\tightlist
\item
  SAF-001: Der Bot muss bei Gefahr sofort anhalten
\item
  SAF-003: Der Bot darf den Tisch nicht verlassen
\item
  SAF-006: Der Bot muss einen Not-Aus unterstützen
\end{itemize}

\paragraph{Beteiligte Komponenten}\label{beteiligte-komponenten-3}

\begin{itemize}
\tightlist
\item
  Sicherheitskomponente
\item
  Bewegungssteuerung
\item
  Aktorsteuerung
\end{itemize}

\paragraph{Voraussetzungen}\label{voraussetzungen-12}

\begin{itemize}
\tightlist
\item
  Roboter befindet sich in Bewegung
\item
  Sicherheitsüberwachung ist aktiv
\end{itemize}

\paragraph{Eingaben}\label{eingaben-12}

Gefahrensignal:

\begin{verbatim}
Tischkante erkannt = TRUE
\end{verbatim}

\paragraph{Ablauf}\label{ablauf-12}

\begin{enumerate}
\def\labelenumi{\arabic{enumi}.}
\tightlist
\item
  Roboter fährt eine definierte Route
\item
  Sensorik meldet kritische Situation
\item
  Sicherheitskomponente erzeugt Stoppsignal
\item
  Bewegungssteuerung beendet Bewegung
\end{enumerate}

\paragraph{Erwartetes Ergebnis}\label{erwartetes-ergebnis-12}

\begin{itemize}
\tightlist
\item
  Bewegung wird sofort beendet
\item
  Aktoren wechseln in sicheren Zustand
\item
  Fehler- oder Warnstatus wird erzeugt
\end{itemize}

\paragraph{Bewertung}\label{bewertung-12}

Bestanden, wenn keine weitere Bewegung nach Aktivierung der
Sicherheitsfunktion erfolgt.

\subsection{Integrationstest: Untersetzerhandling und
Navigation}\label{integrationstest-untersetzerhandling-und-navigation}

\subsubsection{Testfall IT-NAV-CST-001: Navigation zu einem Untersetzer
und
Aufnahme}\label{testfall-it-nav-cst-001-navigation-zu-einem-untersetzer-und-aufnahme}

\paragraph{Ziel}\label{ziel-13}

Überprüfung des Zusammenspiels zwischen Navigation und
Untersetzerhandling.

\paragraph{Zugeordnete
Anforderungen}\label{zugeordnete-anforderungen-13}

\begin{itemize}
\tightlist
\item
  CST-001: Der Bot muss einen Getränkeuntersetzer aufnehmen
\item
  CST-005: Der Bot muss erkennen, ob ein Untersetzer aufgenommen wurde
\item
  NAV-008: Der Bot muss Zielpunkte erreichen
\end{itemize}

\paragraph{Beteiligte Komponenten}\label{beteiligte-komponenten-4}

\begin{itemize}
\tightlist
\item
  Navigation
\item
  Lokalisierung
\item
  Untersetzerhandling
\item
  Systemüberwachung
\end{itemize}

\paragraph{Voraussetzungen}\label{voraussetzungen-13}

\begin{itemize}
\tightlist
\item
  Untersetzerposition ist bekannt
\item
  Roboter befindet sich in Startposition
\end{itemize}

\paragraph{Eingaben}\label{eingaben-13}

Untersetzerposition:

\begin{verbatim}
x = 700 mm
y = 400 mm
\end{verbatim}

\paragraph{Ablauf}\label{ablauf-13}

\begin{enumerate}
\def\labelenumi{\arabic{enumi}.}
\tightlist
\item
  Zielposition des Untersetzers wird übergeben
\item
  Navigation berechnet Route
\item
  Roboter fährt zur Position
\item
  Aufnahmeprozess wird gestartet
\item
  Sensor überprüft Aufnahme
\end{enumerate}

\paragraph{Erwartetes Ergebnis}\label{erwartetes-ergebnis-13}

\begin{itemize}
\tightlist
\item
  Roboter erreicht Untersetzerposition
\item
  Aufnahme wird durchgeführt
\item
  Erfolgreiche Aufnahme wird erkannt
\end{itemize}

\paragraph{Bewertung}\label{bewertung-13}

Bestanden, wenn der vollständige Ablauf ohne Fehler abgeschlossen wird.

\subsection{Integrationstest: Systemüberwachung und
Komponenten}\label{integrationstest-systemuxfcberwachung-und-komponenten}

\subsubsection{Testfall IT-MON-001: Fehlerweitergabe im
Gesamtsystem}\label{testfall-it-mon-001-fehlerweitergabe-im-gesamtsystem}

\paragraph{Ziel}\label{ziel-14}

Überprüfung der Weitergabe und Verarbeitung von Fehlerzuständen.

\paragraph{Zugeordnete
Anforderungen}\label{zugeordnete-anforderungen-14}

\begin{itemize}
\tightlist
\item
  MON-001: Der Bot muss Fehler erkennen
\item
  MON-002: Der Bot muss Fehler protokollieren
\item
  MON-005: Der Bot muss fehlgeschlagene Aufgaben melden
\end{itemize}

\paragraph{Beteiligte Komponenten}\label{beteiligte-komponenten-5}

\begin{itemize}
\tightlist
\item
  Sensorik
\item
  Systemüberwachung
\item
  Sicherheitskomponente
\item
  Benutzerstatusanzeige
\end{itemize}

\paragraph{Voraussetzungen}\label{voraussetzungen-14}

\begin{itemize}
\tightlist
\item
  System befindet sich im Betriebszustand
\end{itemize}

\paragraph{Eingaben}\label{eingaben-14}

Simulierter Fehler:

\begin{verbatim}
Sensorfehler = TRUE
\end{verbatim}

\paragraph{Ablauf}\label{ablauf-14}

\begin{enumerate}
\def\labelenumi{\arabic{enumi}.}
\tightlist
\item
  Sensor meldet Fehler
\item
  Systemüberwachung erkennt Fehler
\item
  Fehler wird gespeichert
\item
  Sicherheitskomponente bewertet Zustand
\item
  Statusmeldung wird erzeugt
\end{enumerate}

\paragraph{Erwartetes Ergebnis}\label{erwartetes-ergebnis-14}

\begin{itemize}
\tightlist
\item
  Fehler wird erkannt
\item
  Fehler wird protokolliert
\item
  System wechselt gegebenenfalls in sicheren Zustand
\item
  Benutzer erhält Statusinformation
\end{itemize}

\paragraph{Bewertung}\label{bewertung-14}

Bestanden, wenn der Fehler vollständig verarbeitet wird.

\subsection{Integrationstest: Hardwareabstraktion und
Simulation}\label{integrationstest-hardwareabstraktion-und-simulation}

\subsubsection{Testfall IT-HW-SIM-001: Austausch realer Hardware durch
Simulation}\label{testfall-it-hw-sim-001-austausch-realer-hardware-durch-simulation}

\paragraph{Ziel}\label{ziel-15}

Überprüfung, ob Softwarekomponenten ohne Anpassung mit simulierten
Hardwarekomponenten betrieben werden können.

\paragraph{Zugeordnete
Anforderungen}\label{zugeordnete-anforderungen-15}

\begin{itemize}
\tightlist
\item
  NFR-004: Software muss simuliert werden können
\item
  NFR-007: Sensoren und Aktoren müssen austauschbar sein
\end{itemize}

\paragraph{Beteiligte Komponenten}\label{beteiligte-komponenten-6}

\begin{itemize}
\tightlist
\item
  Hardwareabstraktion
\item
  Simulationsmodelle
\item
  Robotersoftware
\end{itemize}

\paragraph{Voraussetzungen}\label{voraussetzungen-15}

\begin{itemize}
\tightlist
\item
  Simulationsumgebung ist gestartet
\item
  Virtuelle Sensoren und Aktoren sind verfügbar
\end{itemize}

\paragraph{Ablauf}\label{ablauf-15}

\begin{enumerate}
\def\labelenumi{\arabic{enumi}.}
\tightlist
\item
  Reale Hardwaretreiber werden durch Simulationsmodelle ersetzt
\item
  Software wird gestartet
\item
  Testablauf wird ausgeführt
\item
  Datenfluss wird überprüft
\end{enumerate}

\paragraph{Erwartetes Ergebnis}\label{erwartetes-ergebnis-15}

\begin{itemize}
\tightlist
\item
  Software startet ohne Anpassungen
\item
  Simulierte Sensorwerte werden verarbeitet
\item
  Aktorbefehle werden korrekt erzeugt
\end{itemize}

\paragraph{Bewertung}\label{bewertung-15}

Bestanden, wenn die Software vollständig in der Simulation betrieben
werden kann.

\subsection{Zusammenfassung der
Integrationstests}\label{zusammenfassung-der-integrationstests}

Die Integrationstests überprüfen die Zusammenarbeit der zentralen
Systemkomponenten des Coasterbots.

Die wichtigsten geprüften Integrationen sind:

\begin{itemize}
\tightlist
\item
  Sensorik und Navigation
\item
  Lokalisierung und Navigation
\item
  Navigation und Bewegungssteuerung
\item
  Sicherheitslogik und Aktorsteuerung
\item
  Navigation und Untersetzerhandling
\item
  Systemüberwachung und Fehlerbehandlung
\item
  Hardwareabstraktion und Simulation
\end{itemize}

Durch erfolgreiche Integrationstests wird sichergestellt, dass die
einzelnen Softwaremodule nicht nur isoliert funktionieren, sondern als
zusammenhängendes System die Anforderungen des Pflichtenhefts erfüllen.

\section{Simulationstests}\label{simulationstests-1}

\subsection{Ziel der Simulationstests}\label{ziel-der-simulationstests}

Simulationstests dienen der Überprüfung des Verhaltens des
Coasterbot-Systems innerhalb einer virtuellen Umgebung. Sie ermöglichen
die Validierung von Softwarefunktionen, bevor diese auf der realen
Hardware ausgeführt werden.

Die Simulation bildet dabei die für den Betrieb relevanten Eigenschaften
des Systems nach:

\begin{itemize}
\tightlist
\item
  Tischumgebung
\item
  Roboterbewegung
\item
  Sensorverhalten
\item
  Aktorverhalten
\item
  Hindernisse
\item
  Fehlerzustände
\end{itemize}

Das Ziel besteht darin, nachzuweisen, dass die Software auch unter
kontrollierten und reproduzierbaren Bedingungen die erwarteten
Systemreaktionen zeigt.

Die Simulationstests erfüllen insbesondere folgende Aufgaben:

\begin{itemize}
\tightlist
\item
  Überprüfung autonomer Navigation
\item
  Validierung der Sicherheitsfunktionen
\item
  Prüfung der Fehlerbehandlung
\item
  Nachweis der Reproduzierbarkeit von Fahrmanövern
\item
  Unterstützung automatisierter Tests
\end{itemize}

\subsection{Testumgebung der
Simulation}\label{testumgebung-der-simulation}

Die Simulation basiert auf der im Pflichtenheft beschriebenen
Simulationsarchitektur.

Die Testumgebung besteht aus folgenden Komponenten:

\begin{longtable}[]{@{}ll@{}}
\toprule\noalign{}
Komponente & Aufgabe \\
\midrule\noalign{}
\endhead
\bottomrule\noalign{}
\endlastfoot
Simulationsumgebung & Abbildung der virtuellen Tischumgebung \\
Robotermodell & Simulation der Bewegung und Position \\
Sensormodelle & Erzeugung virtueller Messwerte \\
Aktormodelle & Simulation der Motor- und Mechanikbefehle \\
Teststeuerung & Automatisierte Durchführung von Szenarien \\
Testauswertung & Vergleich von Soll- und Istverhalten \\
\end{longtable}

Die Software des Coasterbots wird dabei möglichst unverändert aus der
realen Umgebung übernommen. Lediglich hardwarenahe Komponenten werden
durch Simulationsmodelle ersetzt.

\subsection{Anforderungen an
Simulationstests}\label{anforderungen-an-simulationstests}

Die Simulationstests basieren auf den im Pflichtenheft definierten
Simulationsanforderungen.

Folgende Eigenschaften müssen erfüllt werden:

\begin{longtable}[]{@{}
  >{\raggedright\arraybackslash}p{(\linewidth - 2\tabcolsep) * \real{0.2535}}
  >{\raggedright\arraybackslash}p{(\linewidth - 2\tabcolsep) * \real{0.7465}}@{}}
\toprule\noalign{}
\begin{minipage}[b]{\linewidth}\raggedright
Anforderung
\end{minipage} & \begin{minipage}[b]{\linewidth}\raggedright
Beschreibung
\end{minipage} \\
\midrule\noalign{}
\endhead
\bottomrule\noalign{}
\endlastfoot
Reproduzierbarkeit & Gleiche Eingaben erzeugen gleiche Ergebnisse \\
Sensorabstraktion & Sensorwerte können simuliert werden \\
Aktorabstraktion & Aktorverhalten kann modelliert werden \\
Fehlerfähigkeit & Fehlerzustände können gezielt erzeugt werden \\
Testunterstützung & Alle definierten Testfälle können ausgeführt
werden \\
Validierung & Simuliertes Verhalten entspricht erwarteten Verhalten \\
\end{longtable}

\subsection{Simulationstest: Grundlegende
Navigation}\label{simulationstest-grundlegende-navigation}

\subsubsection{Testfall ST-SIM-001: Navigation zu einem
Zielpunkt}\label{testfall-st-sim-001-navigation-zu-einem-zielpunkt}

\paragraph{Ziel}\label{ziel-16}

Überprüfung der grundlegenden Navigationsfähigkeit innerhalb der
simulierten Tischumgebung.

\paragraph{Zugeordnete
Anforderungen}\label{zugeordnete-anforderungen-16}

\begin{itemize}
\tightlist
\item
  NAV-006: Der Bot muss seine Position bestimmen können
\item
  NAV-008: Der Bot muss seinen Zielpunkt erreichen
\item
  SIM-001: Fahrmanöver müssen reproduzierbar sein
\end{itemize}

\paragraph{Simulationsumgebung}\label{simulationsumgebung-1}

\begin{itemize}
\tightlist
\item
  Tischgröße: definiert
\item
  Keine Hindernisse vorhanden
\item
  Roboterposition bekannt
\item
  Zielposition bekannt
\end{itemize}

\paragraph{Voraussetzungen}\label{voraussetzungen-16}

\begin{itemize}
\tightlist
\item
  Simulation ist gestartet
\item
  Robotermodell ist initialisiert
\item
  Navigationskomponente ist aktiv
\end{itemize}

\paragraph{Eingaben}\label{eingaben-15}

Startposition:

\begin{verbatim}
x = 0 mm
y = 0 mm
Orientierung = 0°
\end{verbatim}

Zielposition:

\begin{verbatim}
x = 1000 mm
y = 500 mm
\end{verbatim}

\paragraph{Ablauf}\label{ablauf-16}

\begin{enumerate}
\def\labelenumi{\arabic{enumi}.}
\tightlist
\item
  Simulation wird gestartet
\item
  Roboter wird an Startposition platziert
\item
  Zielpunkt wird an Navigation übergeben
\item
  Navigation berechnet Route
\item
  Bewegungssteuerung führt Fahrbefehle aus
\item
  Position wird während der Bewegung überwacht
\end{enumerate}

\paragraph{Erwartetes Ergebnis}\label{erwartetes-ergebnis-16}

\begin{itemize}
\tightlist
\item
  Roboter bewegt sich entlang einer gültigen Route
\item
  Zielpunkt wird erreicht
\item
  Keine unerwarteten Bewegungen treten auf
\item
  Ergebnis ist bei Wiederholung identisch
\end{itemize}

\paragraph{Bewertung}\label{bewertung-16}

Bestanden, wenn der Zielpunkt innerhalb der definierten Genauigkeit
erreicht wird.

\subsection{Simulationstest: Hinderniserkennung und
Ausweichverhalten}\label{simulationstest-hinderniserkennung-und-ausweichverhalten}

\subsubsection{Testfall ST-SIM-002: Dynamische
Hindernisvermeidung}\label{testfall-st-sim-002-dynamische-hindernisvermeidung}

\paragraph{Ziel}\label{ziel-17}

Überprüfung, ob der Roboter Hindernisse erkennt und seine Route
selbstständig anpasst.

\paragraph{Zugeordnete
Anforderungen}\label{zugeordnete-anforderungen-17}

\begin{itemize}
\tightlist
\item
  NAV-001: Hindernisse erkennen
\item
  NAV-002: Hindernissen ausweichen
\item
  NAV-005: Fahrtroute dynamisch anpassen
\item
  SIM-004: Hindernisse müssen berücksichtigt werden
\end{itemize}

\paragraph{Simulationsumgebung}\label{simulationsumgebung-2}

\begin{itemize}
\tightlist
\item
  Tischfläche vorhanden
\item
  Hindernis wird auf geplanter Route platziert
\item
  Zielpunkt erreichbar
\end{itemize}

\paragraph{Voraussetzungen}\label{voraussetzungen-17}

\begin{itemize}
\tightlist
\item
  Navigation ist aktiv
\item
  Hindernissensoren funktionieren
\end{itemize}

\paragraph{Eingaben}\label{eingaben-16}

Hindernis:

\begin{verbatim}
Position = (500 mm, 250 mm)
Größe = 100 mm x 100 mm
\end{verbatim}

\subparagraph{Ablauf}\label{ablauf-17}

\begin{enumerate}
\def\labelenumi{\arabic{enumi}.}
\tightlist
\item
  Navigation startet Fahrt zum Ziel
\item
  Roboter nähert sich Hindernis
\item
  Sensorsimulation erzeugt Hinderniserkennung
\item
  Navigation berechnet neue Route
\item
  Bewegung wird fortgesetzt
\end{enumerate}

\paragraph{Erwartetes Ergebnis}\label{erwartetes-ergebnis-17}

\begin{itemize}
\tightlist
\item
  Hindernis wird erkannt
\item
  Kollision wird verhindert
\item
  Neue Route wird berechnet
\item
  Zielpunkt wird weiterhin erreicht
\end{itemize}

\paragraph{Bewertung}\label{bewertung-17}

Bestanden, wenn der Roboter das Hindernis vollständig umgeht.

\subsection{Simulationstest:
Tischkantenerkennung}\label{simulationstest-tischkantenerkennung}

\subsubsection{Testfall ST-SIM-003: Verhindern des Verlassens der
Tischfläche}\label{testfall-st-sim-003-verhindern-des-verlassens-der-tischfluxe4che}

\paragraph{Ziel}\label{ziel-18}

Überprüfung der Sicherheitsreaktion bei Annäherung an eine Tischkante.

\paragraph{Zugeordnete
Anforderungen}\label{zugeordnete-anforderungen-18}

\begin{itemize}
\tightlist
\item
  NAV-003: Tischende erkennen
\item
  NAV-004: Auf dem Tisch bleiben
\item
  SAF-003: Tisch nicht verlassen
\end{itemize}

\paragraph{Simulationsumgebung}\label{simulationsumgebung-3}

\begin{itemize}
\tightlist
\item
  Tischrand wird modelliert
\item
  Roboter fährt kontrolliert Richtung Rand
\end{itemize}

\paragraph{Voraussetzungen}\label{voraussetzungen-18}

\begin{itemize}
\tightlist
\item
  Tischkantensensor ist aktiv
\item
  Sicherheitslogik ist aktiviert
\end{itemize}

\paragraph{Ablauf}\label{ablauf-18}

\begin{enumerate}
\def\labelenumi{\arabic{enumi}.}
\tightlist
\item
  Roboter bewegt sich Richtung Tischkante
\item
  Simulation erzeugt kritische Abstandswerte
\item
  Sicherheitskomponente verarbeitet Signal
\item
  Bewegungssteuerung erhält Stoppsignal
\end{enumerate}

\paragraph{Erwartetes Ergebnis}\label{erwartetes-ergebnis-18}

\begin{itemize}
\tightlist
\item
  Tischkante wird erkannt
\item
  Roboter stoppt vor dem Rand
\item
  Tischfläche wird nicht verlassen
\end{itemize}

\paragraph{Bewertung}\label{bewertung-18}

Bestanden, wenn der Roboter sicher innerhalb der Tischfläche verbleibt.

\subsection{Simulationstest: Reproduzierbarkeit von
Fahrmanövern}\label{simulationstest-reproduzierbarkeit-von-fahrmanuxf6vern}

\subsubsection{Testfall ST-SIM-004: Wiederholbarkeit identischer
Bewegungsabläufe}\label{testfall-st-sim-004-wiederholbarkeit-identischer-bewegungsabluxe4ufe}

\paragraph{Ziel}\label{ziel-19}

Überprüfung, ob gleiche Eingangsbedingungen zu reproduzierbaren
Ergebnissen führen.

\paragraph{Zugeordnete
Anforderungen}\label{zugeordnete-anforderungen-19}

\begin{itemize}
\tightlist
\item
  SIM-001: Fahrmanöver müssen reproduzierbar sein
\item
  SIM-007: Simulationsverhalten muss erwartbarem Verhalten entsprechen
\end{itemize}

\paragraph{Voraussetzungen}\label{voraussetzungen-19}

\begin{itemize}
\tightlist
\item
  Identisches Szenario kann mehrfach gestartet werden
\end{itemize}

\paragraph{Ablauf}\label{ablauf-19}

\begin{enumerate}
\def\labelenumi{\arabic{enumi}.}
\tightlist
\item
  Simulation wird mit definierten Parametern gestartet
\item
  Navigation wird ausgeführt
\item
  Bewegungsverlauf wird gespeichert
\item
  Simulation wird erneut mit gleichen Parametern gestartet
\item
  Ergebnisse werden verglichen
\end{enumerate}

\paragraph{Erwartetes Ergebnis}\label{erwartetes-ergebnis-19}

\begin{itemize}
\tightlist
\item
  Fahrwege stimmen überein
\item
  Zustandsübergänge erfolgen identisch
\item
  Ergebnisse sind vergleichbar
\end{itemize}

\paragraph{Bewertung}\label{bewertung-19}

Bestanden, wenn die Abweichungen innerhalb definierter Grenzen liegen.

\subsection{Simulationstest:
Sensorausfall}\label{simulationstest-sensorausfall}

\subsubsection{Testfall ST-SIM-005: Reaktion auf fehlerhafte
Sensordaten}\label{testfall-st-sim-005-reaktion-auf-fehlerhafte-sensordaten}

\paragraph{Ziel}\label{ziel-20}

Überprüfung des Systemverhaltens bei einem simulierten Sensorausfall.

\paragraph{Zugeordnete
Anforderungen}\label{zugeordnete-anforderungen-20}

\begin{itemize}
\tightlist
\item
  SAF-002: Bei Sensorfehler sicheren Zustand einnehmen
\item
  MON-001: Fehler erkennen
\item
  MON-002: Fehler protokollieren
\item
  SIM-005: Fehlersituationen simulieren können
\end{itemize}

\paragraph{Voraussetzungen}\label{voraussetzungen-20}

\begin{itemize}
\tightlist
\item
  Roboter befindet sich im Fahrbetrieb
\item
  Fehlersimulation ist aktiviert
\end{itemize}

\paragraph{Eingaben}\label{eingaben-17}

Fehler:

\begin{verbatim}
Abstandssensor = nicht verfügbar
\end{verbatim}

\paragraph{Ablauf}\label{ablauf-20}

\begin{enumerate}
\def\labelenumi{\arabic{enumi}.}
\tightlist
\item
  Simulation startet normalen Betrieb
\item
  Sensorfehler wird ausgelöst
\item
  Systemüberwachung erkennt Fehler
\item
  Sicherheitskomponente bewertet Zustand
\end{enumerate}

\paragraph{Erwartetes Ergebnis}\label{erwartetes-ergebnis-20}

\begin{itemize}
\tightlist
\item
  Fehler wird erkannt
\item
  Fehler wird protokolliert
\item
  Roboter stoppt oder wechselt in sicheren Zustand
\item
  Statusmeldung wird erzeugt
\end{itemize}

\paragraph{Bewertung}\label{bewertung-20}

Bestanden, wenn das System keine unsichere Bewegung fortsetzt.

\subsection{Simulationstest:
Aktorfehler}\label{simulationstest-aktorfehler}

\subsubsection{Testfall ST-SIM-006: Verhalten bei fehlerhafter
Motorsteuerung}\label{testfall-st-sim-006-verhalten-bei-fehlerhafter-motorsteuerung}

\paragraph{Ziel}\label{ziel-21}

Überprüfung der Fehlerbehandlung bei einem Aktorausfall.

\paragraph{Zugeordnete
Anforderungen}\label{zugeordnete-anforderungen-21}

\begin{itemize}
\tightlist
\item
  MON-001: Fehler erkennen
\item
  MON-005: Fehlgeschlagene Aufgaben melden
\item
  SIM-003: Aktoren können simuliert werden
\end{itemize}

\paragraph{Voraussetzungen}\label{voraussetzungen-21}

\begin{itemize}
\tightlist
\item
  Roboter befindet sich in Bewegung
\end{itemize}

\paragraph{Eingaben}\label{eingaben-18}

Fehler:

\begin{verbatim}
Motor links = ausgefallen
\end{verbatim}

\paragraph{Ablauf}\label{ablauf-21}

\begin{enumerate}
\def\labelenumi{\arabic{enumi}.}
\tightlist
\item
  Bewegungsbefehl wird erzeugt
\item
  Simulierter Aktor reagiert fehlerhaft
\item
  Systemüberwachung erkennt Abweichung
\item
  Sicherheitslogik wird aktiviert
\end{enumerate}

\paragraph{Erwartetes Ergebnis}\label{erwartetes-ergebnis-21}

\begin{itemize}
\tightlist
\item
  Aktorfehler wird erkannt
\item
  Bewegung wird kontrolliert beendet
\item
  Fehler wird gemeldet
\end{itemize}

\paragraph{Bewertung}\label{bewertung-21}

Bestanden, wenn das System sicher reagiert.

\subsection{Simulationstest:
Untersetzerhandling}\label{simulationstest-untersetzerhandling}

\subsubsection{Testfall ST-SIM-007: Aufnahme und Platzierung eines
Untersetzers}\label{testfall-st-sim-007-aufnahme-und-platzierung-eines-untersetzers}

\paragraph{Ziel}\label{ziel-22}

Überprüfung des vollständigen Untersetzerprozesses innerhalb der
Simulation.

\paragraph{Zugeordnete
Anforderungen}\label{zugeordnete-anforderungen-22}

\begin{itemize}
\tightlist
\item
  CST-001: Untersetzer aufnehmen
\item
  CST-003: Untersetzer platzieren
\item
  CST-005: Aufnahme erkennen
\item
  SIM-006: Simulation unterstützt Testfälle
\end{itemize}

\paragraph{Voraussetzungen}\label{voraussetzungen-22}

\begin{itemize}
\tightlist
\item
  Untersetzer ist virtuell vorhanden
\item
  Zielposition ist definiert
\end{itemize}

\paragraph{Ablauf}\label{ablauf-22}

\begin{enumerate}
\def\labelenumi{\arabic{enumi}.}
\tightlist
\item
  Roboter navigiert zum Untersetzer
\item
  Aufnahmevorgang wird gestartet
\item
  Zustand wird überprüft
\item
  Roboter fährt zur Zielposition
\item
  Untersetzer wird platziert
\end{enumerate}

\paragraph{Erwartetes Ergebnis}\label{erwartetes-ergebnis-22}

\begin{itemize}
\tightlist
\item
  Untersetzer wird erkannt
\item
  Aufnahme funktioniert
\item
  Transport erfolgt
\item
  Platzierung wird bestätigt
\end{itemize}

\paragraph{Bewertung}\label{bewertung-22}

Bestanden, wenn der komplette Prozess ohne Fehler abgeschlossen wird.

\subsection{Simulationstest: Validierung Simulation gegen erwartetes
Verhalten}\label{simulationstest-validierung-simulation-gegen-erwartetes-verhalten}

\subsubsection{Testfall ST-SIM-008: Vergleich Simulation und
Systemverhalten}\label{testfall-st-sim-008-vergleich-simulation-und-systemverhalten}

\paragraph{Ziel}\label{ziel-23}

Überprüfung, ob die Simulation das erwartete Verhalten des realen
Systems ausreichend genau abbildet.

\paragraph{Zugeordnete
Anforderungen}\label{zugeordnete-anforderungen-23}

\begin{itemize}
\tightlist
\item
  SIM-007: Simulationsverhalten entspricht erwartetem Verhalten
\end{itemize}

\paragraph{Ablauf}\label{ablauf-23}

\begin{enumerate}
\def\labelenumi{\arabic{enumi}.}
\tightlist
\item
  Test wird zunächst in der Simulation durchgeführt
\item
  Ergebnis wird dokumentiert
\item
  Vergleichbare Situation wird am Prototyp durchgeführt
\item
  Ergebnisse werden verglichen
\end{enumerate}

\paragraph{Erwartetes Ergebnis}\label{erwartetes-ergebnis-23}

\begin{itemize}
\tightlist
\item
  Bewegungsabläufe sind vergleichbar
\item
  Zustandsübergänge stimmen überein
\item
  Abweichungen sind erklärbar
\end{itemize}

\paragraph{Bewertung}\label{bewertung-23}

Bestanden, wenn die Simulation als valide Entwicklungsumgebung genutzt
werden kann.

\subsection{Zusammenfassung der
Simulationstests}\label{zusammenfassung-der-simulationstests}

Die Simulationstests überprüfen das Verhalten des Coasterbot-Systems
unter kontrollierten und reproduzierbaren Bedingungen.

Geprüft werden insbesondere:

\begin{itemize}
\tightlist
\item
  autonome Navigation
\item
  Hindernisvermeidung
\item
  Tischkantenerkennung
\item
  Fehlerbehandlung
\item
  Sensor- und Aktormodelle
\item
  Untersetzerhandling
\item
  Validierung der Simulation
\end{itemize}

Durch die Simulation können kritische Situationen sicher getestet und
Entwicklungszyklen verkürzt werden. Gleichzeitig stellt sie sicher, dass
die Softwarearchitektur unabhängig von der realen Hardware überprüfbar
bleibt.

Die erfolgreichen Simulationstests bilden die Grundlage für die
abschließenden Systemtests am realen Prototyp.

\section{Systemtests}\label{systemtests-1}

\subsection{Ziel der Systemtests}\label{ziel-der-systemtests}

Systemtests überprüfen das vollständige Coasterbot-System unter
realitätsnahen Bedingungen. Während Komponenten- und Integrationstests
einzelne Softwaremodule beziehungsweise deren Zusammenspiel betrachten,
wird beim Systemtest der gesamte Prototyp als Gesamtsystem bewertet.

Ziel ist es, nachzuweisen, dass der entwickelte Coasterbot die im
Lasten- und Pflichtenheft definierten Kernfunktionen erfüllt.

Die Systemtests betrachten dabei:

\begin{itemize}
\tightlist
\item
  mechanische Funktionalität
\item
  elektrische Komponenten
\item
  Sensorik
\item
  Aktorik
\item
  Software
\item
  Sicherheitsfunktionen
\item
  Benutzerinteraktion über Statusmeldungen
\end{itemize}

Die Systemtests stellen die abschließende technische Verifikation des
Prototyps dar.

\subsection{Testumgebung}\label{testumgebung-1}

Die Systemtests werden am vollständigen Coasterbot-Prototyp
durchgeführt.

Die Testumgebung umfasst:

\begin{longtable}[]{@{}
  >{\raggedright\arraybackslash}p{(\linewidth - 2\tabcolsep) * \real{0.2209}}
  >{\raggedright\arraybackslash}p{(\linewidth - 2\tabcolsep) * \real{0.7791}}@{}}
\toprule\noalign{}
\begin{minipage}[b]{\linewidth}\raggedright
Komponente
\end{minipage} & \begin{minipage}[b]{\linewidth}\raggedright
Beschreibung
\end{minipage} \\
\midrule\noalign{}
\endhead
\bottomrule\noalign{}
\endlastfoot
Prototyp-Hardware & Vollständig montierter Coasterbot \\
Tischumgebung & Realer Testtisch mit definierter Fläche \\
Hindernisse & Testobjekte zur Simulation realer Situationen \\
Getränkeuntersetzer & Testobjekte für Aufnahme und Platzierung \\
Messmittel & Werkzeuge zur Bewertung von Position, Geschwindigkeit und
Verhalten \\
Testsoftware & Werkzeuge zur Protokollierung und Auswertung \\
\end{longtable}

Die Testumgebung muss möglichst reproduzierbare Bedingungen
gewährleisten.

Dazu gehören:

\begin{itemize}
\tightlist
\item
  definierte Tischgröße
\item
  bekannte Startpositionen
\item
  dokumentierte Hindernispositionen
\item
  definierte Testabläufe
\end{itemize}

\subsection{Systemteststruktur}\label{systemteststruktur}

Jeder Systemtest wird nach folgendem Schema dokumentiert:

\begin{longtable}[]{@{}ll@{}}
\toprule\noalign{}
Feld & Beschreibung \\
\midrule\noalign{}
\endhead
\bottomrule\noalign{}
\endlastfoot
Test-ID & Eindeutige Identifikation \\
Testname & Bezeichnung des Tests \\
Ziel & Zu prüfende Systemfunktion \\
Anforderungen & Zugeordnete Anforderungen \\
Voraussetzungen & Erforderliche Bedingungen \\
Testaufbau & Beschreibung der Umgebung \\
Ablauf & Durchführung \\
Erwartetes Ergebnis & Sollverhalten \\
Bewertung & Testergebnis \\
\end{longtable}

\subsection{Systemtest: Inbetriebnahme}\label{systemtest-inbetriebnahme}

\subsubsection{Testfall SYS-001: Systemstart und
Initialisierung}\label{testfall-sys-001-systemstart-und-initialisierung}

\paragraph{Ziel}\label{ziel-24}

Überprüfung, ob der Coasterbot nach dem Einschalten alle notwendigen
Komponenten korrekt initialisiert.

\paragraph{Zugeordnete
Anforderungen}\label{zugeordnete-anforderungen-24}

\begin{itemize}
\tightlist
\item
  MON-003: Der Bot muss seinen Betriebszustand anzeigen
\item
  NFR-001: Der Bot muss modular aufgebaut sein
\item
  NFR-005: Die Software muss testbar sein
\end{itemize}

\paragraph{Voraussetzungen}\label{voraussetzungen-23}

\begin{itemize}
\tightlist
\item
  Hardware ist vollständig montiert
\item
  Energieversorgung ist vorhanden
\item
  Software ist installiert
\end{itemize}

\paragraph{Testaufbau}\label{testaufbau}

Der Coasterbot befindet sich ausgeschaltet auf einer ebenen Tischfläche.

\paragraph{Ablauf}\label{ablauf-24}

\begin{enumerate}
\def\labelenumi{\arabic{enumi}.}
\tightlist
\item
  Energieversorgung wird aktiviert
\item
  Steuerungssystem startet
\item
  Hardwarekomponenten werden initialisiert
\item
  Sensoren werden geprüft
\item
  Betriebszustand wird ausgegeben
\end{enumerate}

\paragraph{Erwartetes Ergebnis}\label{erwartetes-ergebnis-24}

\begin{itemize}
\tightlist
\item
  Alle Komponenten starten erfolgreich
\item
  Keine kritischen Fehler werden erkannt
\item
  Systemstatus zeigt Betriebsbereitschaft
\end{itemize}

\paragraph{Bewertung}\label{bewertung-24}

Bestanden, wenn der Roboter vollständig betriebsbereit ist.

\subsection{Systemtest: Navigation auf
Tischfläche}\label{systemtest-navigation-auf-tischfluxe4che}

\subsubsection{Testfall SYS-002: Autonome Bewegung auf dem
Tisch}\label{testfall-sys-002-autonome-bewegung-auf-dem-tisch}

\paragraph{Ziel}\label{ziel-25}

Überprüfung der grundlegenden autonomen Bewegungsfähigkeit des
Prototyps.

\paragraph{Zugeordnete
Anforderungen}\label{zugeordnete-anforderungen-25}

\begin{itemize}
\tightlist
\item
  NAV-006: Position bestimmen
\item
  NAV-007: Orientierung bestimmen
\item
  NAV-008: Zielpunkt erreichen
\item
  NAV-009: Kontrollierte Bewegung
\end{itemize}

\paragraph{Voraussetzungen}\label{voraussetzungen-24}

\begin{itemize}
\tightlist
\item
  System ist gestartet
\item
  Tischfläche ist frei
\item
  Start- und Zielposition sind definiert
\end{itemize}

\paragraph{Testaufbau}\label{testaufbau-1}

Der Roboter wird an einer definierten Startposition platziert.

Beispiel:

Start:

\[
(0,0)
\]

Ziel:

\[
(1000,500)
\]

\paragraph{Ablauf}\label{ablauf-25}

\begin{enumerate}
\def\labelenumi{\arabic{enumi}.}
\tightlist
\item
  Zielposition wird vorgegeben
\item
  Navigation wird aktiviert
\item
  Roboter fährt selbstständig zum Ziel
\item
  Position wird aufgezeichnet
\end{enumerate}

\paragraph{Erwartetes Ergebnis}\label{erwartetes-ergebnis-25}

\begin{itemize}
\tightlist
\item
  Roboter bewegt sich selbstständig
\item
  Fahrbewegung bleibt kontrolliert
\item
  Zielpunkt wird erreicht
\item
  Keine ungewollten Bewegungen treten auf
\end{itemize}

\paragraph{Bewertung}\label{bewertung-25}

Bestanden, wenn die Zielposition innerhalb der definierten Genauigkeit
erreicht wird.

\subsection{Systemtest:
Hindernisvermeidung}\label{systemtest-hindernisvermeidung}

\subsubsection{Testfall SYS-003: Erkennung und Umfahrung eines
Hindernisses}\label{testfall-sys-003-erkennung-und-umfahrung-eines-hindernisses}

\paragraph{Ziel}\label{ziel-26}

Überprüfung der Fähigkeit des Roboters, Hindernisse im realen Betrieb zu
erkennen und zu umgehen.

\paragraph{Zugeordnete
Anforderungen}\label{zugeordnete-anforderungen-26}

\begin{itemize}
\tightlist
\item
  NAV-001: Hindernisse erkennen
\item
  NAV-002: Hindernissen ausweichen
\item
  NAV-005: Fahrroute dynamisch anpassen
\item
  SAF-005: Kollisionen vermeiden
\end{itemize}

\paragraph{Voraussetzungen}\label{voraussetzungen-25}

\begin{itemize}
\tightlist
\item
  Navigation ist aktiv
\item
  Hindernis befindet sich auf der geplanten Route
\end{itemize}

\paragraph{Testaufbau}\label{testaufbau-2}

Ein Hindernis wird auf der Tischfläche platziert.

Beispiel:

\begin{itemize}
\tightlist
\item
  Becher
\item
  Gegenstand
\item
  Tischdekoration
\end{itemize}

\paragraph{Ablauf}\label{ablauf-26}

\begin{enumerate}
\def\labelenumi{\arabic{enumi}.}
\tightlist
\item
  Roboter startet Navigation
\item
  Hindernis wird erkannt
\item
  Navigation berechnet alternative Route
\item
  Roboter fährt um Hindernis herum
\end{enumerate}

\paragraph{Erwartetes Ergebnis}\label{erwartetes-ergebnis-26}

\begin{itemize}
\tightlist
\item
  Hindernis wird erkannt
\item
  Keine Kollision entsteht
\item
  Route wird angepasst
\item
  Ziel wird weiterhin erreicht
\end{itemize}

\paragraph{Bewertung}\label{bewertung-26}

Bestanden, wenn das Hindernis sicher umfahren wird.

\subsection{Systemtest:
Tischkantenerkennung}\label{systemtest-tischkantenerkennung}

\subsubsection{Testfall SYS-004: Schutz vor Herunterfallen vom
Tisch}\label{testfall-sys-004-schutz-vor-herunterfallen-vom-tisch}

\paragraph{Ziel}\label{ziel-27}

Überprüfung der Sicherheitsfunktion zur Vermeidung eines Absturzes.

\paragraph{Zugeordnete
Anforderungen}\label{zugeordnete-anforderungen-27}

\begin{itemize}
\tightlist
\item
  NAV-003: Tischende erkennen
\item
  NAV-004: Auf dem Tisch bleiben
\item
  SAF-003: Tisch nicht verlassen
\end{itemize}

\paragraph{Voraussetzungen}\label{voraussetzungen-26}

\begin{itemize}
\tightlist
\item
  Tischkantenerkennung ist aktiviert
\item
  Roboter bewegt sich kontrolliert
\end{itemize}

\paragraph{Testaufbau}\label{testaufbau-3}

Der Roboter wird auf eine Position nahe der Tischkante gefahren.

\paragraph{Ablauf}\label{ablauf-27}

\begin{enumerate}
\def\labelenumi{\arabic{enumi}.}
\tightlist
\item
  Roboter fährt Richtung Tischrand
\item
  Sensorik erkennt Annäherung
\item
  Sicherheitslogik bewertet Situation
\item
  Bewegungssteuerung stoppt
\end{enumerate}

\paragraph{Erwartetes Ergebnis}\label{erwartetes-ergebnis-27}

\begin{itemize}
\tightlist
\item
  Tischkante wird erkannt
\item
  Roboter hält ausreichend Abstand
\item
  Kein Teil des Roboters verlässt die Tischfläche
\end{itemize}

\paragraph{Bewertung}\label{bewertung-27}

Bestanden, wenn der Roboter sicher auf dem Tisch verbleibt.

\subsection{Systemtest:
Untersetzeraufnahme}\label{systemtest-untersetzeraufnahme}

\subsubsection{Testfall SYS-005: Aufnahme eines
Getränkeuntersetzers}\label{testfall-sys-005-aufnahme-eines-getruxe4nkeuntersetzers}

\paragraph{Ziel}\label{ziel-28}

Überprüfung der mechanischen und softwareseitigen Funktion zur Aufnahme
eines Untersetzers.

\paragraph{Zugeordnete
Anforderungen}\label{zugeordnete-anforderungen-28}

\begin{itemize}
\tightlist
\item
  CST-001: Untersetzer aufnehmen
\item
  CST-005: Erfolgreiche Aufnahme erkennen
\end{itemize}

\paragraph{Voraussetzungen}\label{voraussetzungen-27}

\begin{itemize}
\tightlist
\item
  Untersetzer befindet sich an definierter Position
\item
  Aufnahmemechanismus funktioniert
\end{itemize}

\paragraph{Testaufbau}\label{testaufbau-4}

Ein Getränkeuntersetzer wird auf der Tischfläche platziert.

\paragraph{Ablauf}\label{ablauf-28}

\begin{enumerate}
\def\labelenumi{\arabic{enumi}.}
\tightlist
\item
  Roboter navigiert zur Untersetzerposition
\item
  Aufnahmevorgang wird gestartet
\item
  Mechanismus nimmt Untersetzer auf
\item
  Sensorik überprüft Zustand
\end{enumerate}

\paragraph{Erwartetes Ergebnis}\label{erwartetes-ergebnis-28}

\begin{itemize}
\tightlist
\item
  Untersetzer wird aufgenommen
\item
  Aufnahme wird erkannt
\item
  Systemstatus wird aktualisiert
\end{itemize}

\paragraph{Bewertung}\label{bewertung-28}

Bestanden, wenn der Untersetzer sicher aufgenommen wurde.

\subsection{Systemtest:
Untersetzerplatzierung}\label{systemtest-untersetzerplatzierung}

\subsubsection{Testfall SYS-006: Platzierung eines
Untersetzers}\label{testfall-sys-006-platzierung-eines-untersetzers}

\paragraph{Ziel}\label{ziel-29}

Überprüfung der präzisen Ablage eines Untersetzers.

\paragraph{Zugeordnete
Anforderungen}\label{zugeordnete-anforderungen-29}

\begin{itemize}
\tightlist
\item
  CST-003: Untersetzer präzise platzieren
\item
  CST-006: Erfolgreiche Platzierung erkennen
\end{itemize}

\paragraph{Voraussetzungen}\label{voraussetzungen-28}

\begin{itemize}
\tightlist
\item
  Untersetzer wurde aufgenommen
\item
  Zielposition ist bekannt
\end{itemize}

\paragraph{Testaufbau}\label{testaufbau-5}

Eine definierte Ablageposition wird markiert.

\paragraph{Ablauf}\label{ablauf-29}

\begin{enumerate}
\def\labelenumi{\arabic{enumi}.}
\tightlist
\item
  Roboter fährt zur Zielposition
\item
  Platzierungsmechanismus wird aktiviert
\item
  Positionierung wird überprüft
\end{enumerate}

\paragraph{Erwartetes Ergebnis}\label{erwartetes-ergebnis-29}

\begin{itemize}
\tightlist
\item
  Untersetzer wird abgelegt
\item
  Position liegt innerhalb der definierten Toleranz
\item
  Erfolgreiche Ablage wird erkannt
\end{itemize}

\paragraph{Bewertung}\label{bewertung-29}

Bestanden, wenn der Untersetzer korrekt positioniert wurde.

\subsection{Systemtest:
Sicherheitsstopp}\label{systemtest-sicherheitsstopp}

\subsubsection{Testfall SYS-007: Aktivierung des
Not-Aus}\label{testfall-sys-007-aktivierung-des-not-aus}

\paragraph{Ziel}\label{ziel-30}

Überprüfung der Sicherheitsreaktion bei manueller oder automatisierter
Notabschaltung.

\paragraph{Zugeordnete
Anforderungen}\label{zugeordnete-anforderungen-30}

\begin{itemize}
\tightlist
\item
  SAF-001: Bei Gefahr sofort anhalten
\item
  SAF-006: Not-Aus unterstützen
\end{itemize}

\paragraph{Voraussetzungen}\label{voraussetzungen-29}

\begin{itemize}
\tightlist
\item
  Roboter befindet sich in Bewegung
\end{itemize}

\paragraph{Testaufbau}\label{testaufbau-6}

Not-Aus-Funktion wird während der Bewegung ausgelöst.

\paragraph{Ablauf}\label{ablauf-30}

\begin{enumerate}
\def\labelenumi{\arabic{enumi}.}
\tightlist
\item
  Roboter fährt
\item
  Not-Aus wird aktiviert
\item
  Sicherheitslogik verarbeitet Signal
\item
  Bewegungssteuerung stoppt
\end{enumerate}

\paragraph{Erwartetes Ergebnis}\label{erwartetes-ergebnis-30}

\begin{itemize}
\tightlist
\item
  Bewegung wird unverzüglich beendet
\item
  Motoren wechseln in sicheren Zustand
\item
  Systemstatus zeigt Sicherheitsstopp
\end{itemize}

\paragraph{Bewertung}\label{bewertung-30}

Bestanden, wenn der Roboter sicher stoppt.

\subsection{Systemtest:
Fehlerbehandlung}\label{systemtest-fehlerbehandlung}

\subsubsection{Testfall SYS-008: Verhalten bei
Systemfehler}\label{testfall-sys-008-verhalten-bei-systemfehler}

\paragraph{Ziel}\label{ziel-31}

Überprüfung der Reaktion auf einen realitätsnahen Fehlerzustand.

\paragraph{Zugeordnete
Anforderungen}\label{zugeordnete-anforderungen-31}

\begin{itemize}
\tightlist
\item
  MON-001: Fehler erkennen
\item
  MON-002: Fehler protokollieren
\item
  MON-005: Fehlgeschlagene Aufgaben melden
\end{itemize}

\paragraph{Voraussetzungen}\label{voraussetzungen-30}

\begin{itemize}
\tightlist
\item
  System befindet sich im Betriebszustand
\end{itemize}

\paragraph{Testaufbau}\label{testaufbau-7}

Ein Fehler wird gezielt erzeugt.

Beispiele:

\begin{itemize}
\tightlist
\item
  Sensor wird getrennt
\item
  Aktor reagiert nicht
\item
  ungültiger Systemzustand
\end{itemize}

\paragraph{Ablauf}\label{ablauf-31}

\begin{enumerate}
\def\labelenumi{\arabic{enumi}.}
\tightlist
\item
  Fehler wird erzeugt
\item
  Systemüberwachung erkennt Fehler
\item
  Fehler wird protokolliert
\item
  Sicherheitsreaktion wird ausgeführt
\end{enumerate}

\paragraph{Erwartetes Ergebnis}\label{erwartetes-ergebnis-31}

\begin{itemize}
\tightlist
\item
  Fehler wird erkannt
\item
  Fehler wird dokumentiert
\item
  System wechselt in definierten Zustand
\end{itemize}

\paragraph{Bewertung}\label{bewertung-31}

Bestanden, wenn der Fehler sicher verarbeitet wird.

\subsection{Systemtest:
Energieversorgung}\label{systemtest-energieversorgung}

\subsubsection{Testfall SYS-009: Überwachung des
Akkuzustands}\label{testfall-sys-009-uxfcberwachung-des-akkuzustands}

\paragraph{Ziel}\label{ziel-32}

Überprüfung der Überwachung der Energieversorgung.

\paragraph{Zugeordnete
Anforderungen}\label{zugeordnete-anforderungen-32}

\begin{itemize}
\tightlist
\item
  Energieversorgung-001: Akkuladestand überwachen
\item
  Energieversorgung-002: Niedrigen Akkustand erkennen
\item
  Energieversorgung-003: Aufgabe kontrolliert abbrechen
\end{itemize}

\paragraph{Voraussetzungen}\label{voraussetzungen-31}

\begin{itemize}
\tightlist
\item
  Akkuüberwachung ist aktiv
\end{itemize}

\paragraph{Testaufbau}\label{testaufbau-8}

Akkustand wird schrittweise reduziert.

\paragraph{Ablauf}\label{ablauf-32}

\begin{enumerate}
\def\labelenumi{\arabic{enumi}.}
\tightlist
\item
  Roboter arbeitet im Normalbetrieb
\item
  Akkustand unterschreitet Grenzwert
\item
  Energieverwaltung reagiert
\end{enumerate}

\paragraph{Erwartetes Ergebnis}\label{erwartetes-ergebnis-32}

\begin{itemize}
\tightlist
\item
  Niedriger Akkustand wird erkannt
\item
  Statusmeldung wird erzeugt
\item
  Aufgabe wird kontrolliert beendet
\end{itemize}

\paragraph{Bewertung}\label{bewertung-32}

Bestanden, wenn keine unkontrollierte Abschaltung erfolgt.

\subsection{Systemtest: Gesamtablauf}\label{systemtest-gesamtablauf}

\subsubsection{Testfall SYS-010: Vollständiger
Coasterbot-Arbeitsablauf}\label{testfall-sys-010-vollstuxe4ndiger-coasterbot-arbeitsablauf}

\paragraph{Ziel}\label{ziel-33}

Überprüfung des vollständigen Kernprozesses.

\paragraph{Zugeordnete
Anforderungen}\label{zugeordnete-anforderungen-33}

\begin{itemize}
\tightlist
\item
  Navigation
\item
  Lokalisierung
\item
  Untersetzerhandling
\item
  Systemüberwachung
\end{itemize}

\paragraph{Ablauf}\label{ablauf-33}

\begin{enumerate}
\def\labelenumi{\arabic{enumi}.}
\tightlist
\item
  Roboter startet
\item
  Zielposition eines Kunden wird definiert
\item
  Roboter navigiert zur Position
\item
  Untersetzer wird aufgenommen
\item
  Roboter fährt zur Zielposition
\item
  Untersetzer wird platziert
\item
  Aufgabe wird abgeschlossen
\end{enumerate}

\paragraph{Erwartetes Ergebnis}\label{erwartetes-ergebnis-33}

Der vollständige Ablauf wird ohne Fehler durchgeführt.

\paragraph{Bewertung}\label{bewertung-33}

Bestanden, wenn alle Teilfunktionen erfolgreich ausgeführt wurden.

\subsection{Zusammenfassung der
Systemtests}\label{zusammenfassung-der-systemtests}

Die Systemtests überprüfen den Coasterbot als vollständiges
Gesamtsystem.

Dabei werden die zentralen Funktionen des Prototyps bewertet:

\begin{itemize}
\tightlist
\item
  sichere Inbetriebnahme
\item
  autonome Navigation
\item
  Hindernisvermeidung
\item
  Tischkantenerkennung
\item
  Untersetzerhandling
\item
  Sicherheitsfunktionen
\item
  Fehlerbehandlung
\item
  Energieüberwachung
\end{itemize}

Die erfolgreichen Systemtests bestätigen, dass die entwickelten
Hardware- und Softwarekomponenten gemeinsam die Anforderungen des
Pflichtenhefts erfüllen.

Damit bilden die Systemtests die abschließende technische Grundlage für
die Bewertung des Prototyps.

\section{Testergebnisse und
Bewertung}\label{testergebnisse-und-bewertung}

\subsection{Ziel der Testauswertung}\label{ziel-der-testauswertung}

Die Testauswertung dient der systematischen Bewertung der durchgeführten
Tests und der Ermittlung des Entwicklungsstandes des
Coasterbot-Prototyps.

Ziel ist es, festzustellen, in welchem Umfang die Anforderungen aus dem
Lasten- und Pflichtenheft erfüllt werden. Die Ergebnisse der
Testdurchführung werden dokumentiert, analysiert und hinsichtlich ihrer
Bedeutung für das Gesamtsystem bewertet.

Die Auswertung umfasst:

\begin{itemize}
\tightlist
\item
  Bewertung einzelner Testfälle
\item
  Analyse aufgetretener Fehler
\item
  Ermittlung der Testabdeckung
\item
  Bewertung der Systemqualität
\item
  Überprüfung der Definition of Done
\end{itemize}

\subsection{Dokumentation der
Testergebnisse}\label{dokumentation-der-testergebnisse}

Jeder durchgeführte Test wird mit seinem Ergebnis dokumentiert.

Die Testergebnisse werden nach folgendem Schema erfasst:

\begin{longtable}[]{@{}ll@{}}
\toprule\noalign{}
Feld & Beschreibung \\
\midrule\noalign{}
\endhead
\bottomrule\noalign{}
\endlastfoot
Test-ID & Identifikation des Testfalls \\
Testname & Bezeichnung des Tests \\
Datum & Zeitpunkt der Durchführung \\
Testumgebung & Simulation oder Prototyp \\
Ergebnis & Bestanden, Fehlgeschlagen oder Nicht ausführbar \\
Bemerkungen & Zusätzliche Informationen \\
Fehler-ID & Referenz auf Fehlerbeschreibung \\
\end{longtable}

\subsection{Bewertungskriterien}\label{bewertungskriterien}

Die Bewertung der Testfälle erfolgt anhand definierter Kriterien.

\subsubsection{Bestanden}\label{bestanden}

Ein Test gilt als bestanden, wenn:

\begin{itemize}
\tightlist
\item
  das erwartete Ergebnis vollständig erreicht wurde
\item
  keine sicherheitskritischen Fehler auftreten
\item
  das Verhalten reproduzierbar ist
\end{itemize}

\subsubsection{Fehlgeschlagen}\label{fehlgeschlagen}

Ein Test gilt als fehlgeschlagen, wenn:

\begin{itemize}
\tightlist
\item
  die erwartete Funktion nicht erfüllt wird
\item
  unerwartetes Verhalten auftritt
\item
  Sicherheitsanforderungen verletzt werden
\end{itemize}

Fehlgeschlagene Tests werden analysiert und führen gegebenenfalls zu
einer Anpassung der Implementierung.

\subsubsection{Nicht ausführbar}\label{nicht-ausfuxfchrbar}

Ein Test wird als nicht ausführbar bewertet, wenn:

\begin{itemize}
\tightlist
\item
  Voraussetzungen nicht erfüllt sind
\item
  benötigte Hardware nicht verfügbar ist
\item
  Abhängigkeiten fehlen
\end{itemize}

Nicht ausführbare Tests müssen begründet und gegebenenfalls nachgeholt
werden.

\subsubsection{Testabdeckung}\label{testabdeckung}

Die Testabdeckung beschreibt, welcher Anteil der Anforderungen durch
Testfälle überprüft wurde.

Die Bewertung erfolgt anhand der Zuordnung zwischen:

\begin{itemize}
\tightlist
\item
  Lastenheft-Anforderungen
\item
  Pflichtenheft-Komponenten
\item
  Testfällen
\end{itemize}

Eine vollständige Abdeckung bedeutet, dass jede relevante Anforderung
mindestens einem Testfall zugeordnet wurde.

Die Testabdeckung wird in folgende Bereiche unterteilt:

\begin{longtable}[]{@{}ll@{}}
\toprule\noalign{}
Bereich & Bewertung \\
\midrule\noalign{}
\endhead
\bottomrule\noalign{}
\endlastfoot
Navigation und Lokalisierung & Vollständig getestet \\
Sicherheit & Vollständig getestet \\
Untersetzerhandling & Vollständig getestet \\
Systemüberwachung & Vollständig getestet \\
Simulation & Vollständig getestet \\
Erweiterungsfunktionen & Nicht vollständig getestet \\
\end{longtable}

\subsection{Auswertung der
Kernfunktionen}\label{auswertung-der-kernfunktionen}

\subsubsection{Navigation und
Lokalisierung}\label{navigation-und-lokalisierung}

Die Navigation stellt eine zentrale Funktion des Coasterbots dar.

Die Tests überprüfen:

\begin{itemize}
\tightlist
\item
  Positionsbestimmung
\item
  Orientierung
\item
  Zielanfahrt
\item
  Hindernisvermeidung
\item
  Anpassung der Fahrroute
\end{itemize}

Die erfolgreiche Durchführung der Navigations- und Simulationstests
bestätigt die grundsätzliche Fähigkeit des Roboters, sich autonom auf
einer Tischfläche zu bewegen.

\subsubsection{Sicherheitsfunktionen}\label{sicherheitsfunktionen}

Die Sicherheitsfunktionen besitzen aufgrund der möglichen Auswirkungen
eine hohe Priorität.

Geprüft wurden:

\begin{itemize}
\tightlist
\item
  Tischkantenerkennung
\item
  Kollisionsvermeidung
\item
  Not-Aus
\item
  sichere Zustände bei Fehlern
\end{itemize}

Die Tests bestätigen, dass kritische Situationen erkannt und geeignete
Schutzmaßnahmen eingeleitet werden.

\subsubsection{Getränkeuntersetzerhandling}\label{getruxe4nkeuntersetzerhandling}

Das Untersetzerhandling bildet die Kernfunktion des Prototyps.

Die Tests bewerten:

\begin{itemize}
\tightlist
\item
  Erkennung eines Untersetzers
\item
  Aufnahme
\item
  Transport
\item
  Positionierung
\item
  Erkennung erfolgreicher Aktionen
\end{itemize}

Die Ergebnisse zeigen, ob die mechanischen und softwareseitigen
Komponenten ausreichend zusammenarbeiten.

\subsubsection{Systemüberwachung}\label{systemuxfcberwachung}

Die Systemüberwachung stellt sicher, dass Fehlerzustände erkannt und
verarbeitet werden.

Bewertet werden:

\begin{itemize}
\tightlist
\item
  Fehlererkennung
\item
  Fehlerprotokollierung
\item
  Statusanzeige
\item
  Erkennung fehlgeschlagener Aufgaben
\end{itemize}

Eine funktionierende Systemüberwachung ist Voraussetzung für einen
zuverlässigen Betrieb.

\subsection{Fehleranalyse}\label{fehleranalyse}

Während der Testdurchführung erkannte Fehler werden nach ihrer
Kritikalität bewertet.

Die Bewertung erfolgt anhand folgender Kategorien:

\begin{longtable}[]{@{}ll@{}}
\toprule\noalign{}
Kategorie & Beschreibung \\
\midrule\noalign{}
\endhead
\bottomrule\noalign{}
\endlastfoot
Kritischer Fehler & Beeinträchtigt Sicherheit oder Kernfunktion \\
Schwerer Fehler & Verhindert wesentliche Funktionen \\
Mittlerer Fehler & Beeinträchtigt Komfort oder Zuverlässigkeit \\
Geringer Fehler & Keine wesentliche Einschränkung \\
\end{longtable}

\subsubsection{Umgang mit kritischen
Fehlern}\label{umgang-mit-kritischen-fehlern}

Kritische Fehler besitzen höchste Priorität und müssen vor Abschluss der
Projektarbeit behoben oder bewertet werden.

Beispiele:

\begin{itemize}
\tightlist
\item
  Verlust der Tischkantenerkennung
\item
  Unkontrollierte Bewegung
\item
  Fehlfunktion des Not-Aus
\item
  Gefährliche Kollision
\end{itemize}

Ein Prototyp darf nicht als erfolgreich bewertet werden, wenn kritische
Sicherheitsanforderungen nicht erfüllt sind.

\subsubsection{Umgang mit funktionalen
Fehlern}\label{umgang-mit-funktionalen-fehlern}

Funktionale Fehler werden hinsichtlich ihrer Auswirkungen bewertet.

Beispiele:

\begin{itemize}
\tightlist
\item
  Ungenaue Positionierung
\item
  Verzögerte Statusmeldung
\item
  Fehlgeschlagene Untersetzeraufnahme
\end{itemize}

Je nach Bedeutung können diese Fehler behoben oder als bekannte
Einschränkung dokumentiert werden.

\subsection{Bewertung der Definition of
Done}\label{bewertung-der-definition-of-done}

Die Definition of Done des Projekts umfasst mehrere technische Ziele.
Diese werden anhand der Testergebnisse bewertet.

\begin{longtable}[]{@{}
  >{\raggedright\arraybackslash}p{(\linewidth - 2\tabcolsep) * \real{0.4045}}
  >{\raggedright\arraybackslash}p{(\linewidth - 2\tabcolsep) * \real{0.5955}}@{}}
\toprule\noalign{}
\begin{minipage}[b]{\linewidth}\raggedright
Definition of Done
\end{minipage} & \begin{minipage}[b]{\linewidth}\raggedright
Bewertungsmethode
\end{minipage} \\
\midrule\noalign{}
\endhead
\bottomrule\noalign{}
\endlastfoot
STL-Dateien und Schaltpläne erstellt & Prüfung der
Fertigungsunterlagen \\
Softwarearchitektur ausgearbeitet & Prüfung der Architektur- und
Komponentendokumentation \\
Simulation erstellt und validiert & Durchführung der Simulationstests \\
Tests identifiziert und definiert & Prüfung der Testspezifikation \\
\end{longtable}

\subsubsection{Fertigungsunterlagen}\label{fertigungsunterlagen}

Die Fertigungsunterlagen gelten als erfüllt, wenn:

\begin{itemize}
\tightlist
\item
  mechanische Bauteile dokumentiert sind
\item
  STL-Dateien vorhanden sind
\item
  Schaltpläne erstellt wurden
\end{itemize}

\subsubsection{Softwarearchitektur}\label{softwarearchitektur}

Die Softwarearchitektur gilt als erfüllt, wenn:

\begin{itemize}
\tightlist
\item
  Komponenten definiert sind
\item
  Schnittstellen dokumentiert sind
\item
  Hardwareabstraktion umgesetzt wurde
\item
  Erweiterbarkeit gewährleistet ist
\end{itemize}

\subsubsection{Simulation}\label{simulation-1}

Die Simulation gilt als erfüllt, wenn:

\begin{itemize}
\tightlist
\item
  Fahrmanöver reproduzierbar sind
\item
  Sensoren modelliert werden können
\item
  Aktoren simuliert werden können
\item
  Fehlerfälle getestet werden können
\end{itemize}

\subsubsection{Testspezifikation}\label{testspezifikation}

Die Testdefinition gilt als erfüllt, wenn:

\begin{itemize}
\tightlist
\item
  Anforderungen Testfällen zugeordnet wurden
\item
  Testabläufe beschrieben sind
\item
  Bewertungskriterien definiert wurden
\end{itemize}

\subsection{Bewertung der
Projektziele}\label{bewertung-der-projektziele}

Die Testergebnisse werden abschließend hinsichtlich der ursprünglichen
Projektziele bewertet.

Das Hauptziel war die Entwicklung eines funktionsfähigen Prototyps eines
Coasterbots.

Dabei standen folgende Funktionen im Mittelpunkt:

\begin{itemize}
\tightlist
\item
  autonome Navigation auf einem Tisch
\item
  Erkennung und Vermeidung von Hindernissen
\item
  sichere Bewegung
\item
  Aufnahme und Platzierung von Getränkeuntersetzern
\item
  modulare Softwarearchitektur
\item
  simulationsgestützte Entwicklung
\end{itemize}

Die Testergebnisse zeigen, welche Funktionen erfolgreich umgesetzt
wurden und welche Bereiche noch Entwicklungspotenzial besitzen.

\subsubsection{Grenzen der Bewertung}\label{grenzen-der-bewertung}

Die Bewertung berücksichtigt den Projektumfang und die definierten
Abgrenzungen.

Nicht Bestandteil der Bewertung sind:

\begin{itemize}
\tightlist
\item
  automatischer Getränketransport
\item
  Bestellung und Bezahlung
\item
  Reinigung
\item
  Netzwerkfunktionen
\item
  Schwarmverhalten mehrerer Roboter
\item
  EMV-Betrachtungen
\end{itemize}

Diese Funktionen können in zukünftigen Entwicklungsstufen betrachtet
werden.

\section{Zusammenfassung}\label{zusammenfassung-2}

Die Testauswertung stellt den erreichten Entwicklungsstand des
Coasterbot-Prototyps dar.

Durch die Kombination aus Komponenten-, Integrations-, Simulations- und
Systemtests konnte die Funktionsfähigkeit der entwickelten Lösung
systematisch überprüft werden.

Die Testergebnisse ermöglichen eine objektive Bewertung der Umsetzung
und zeigen gleichzeitig mögliche zukünftige Erweiterungen und
Optimierungspotenziale auf.

Damit bildet dieses Kapitel die Grundlage für die abschließende
Bewertung des Projektergebnisses.

\subsection{Zusammenfassung und
Ausblick}\label{zusammenfassung-und-ausblick}

\subsubsection{Zusammenfassung der
Testspezifikation}\label{zusammenfassung-der-testspezifikation}

Diese Testspezifikation beschreibt die systematische Überprüfung des im
Rahmen der Projektarbeit entwickelten Coasterbot-Prototyps.

Ziel der Testdurchführung war es, die im Lasten- und Pflichtenheft
definierten Anforderungen nachvollziehbar zu überprüfen und den
Entwicklungsstand des Prototyps zu bewerten.

Hierzu wurden verschiedene Testebenen betrachtet:

\begin{itemize}
\tightlist
\item
  Komponententests zur Prüfung einzelner Softwaremodule
\item
  Integrationstests zur Überprüfung des Zusammenspiels mehrerer
  Komponenten
\item
  Simulationstests zur Validierung des Systemverhaltens unter
  kontrollierten Bedingungen
\item
  Systemtests zur Bewertung des vollständigen Prototyps
\end{itemize}

Durch diese strukturierte Vorgehensweise konnten sowohl einzelne
Funktionen als auch das Gesamtsystem betrachtet werden.

\subsubsection{Bewertung der erreichten
Projektziele}\label{bewertung-der-erreichten-projektziele}

Das zentrale Ziel der Projektarbeit war die Entwicklung eines Prototyps
eines Coasterbots, der sich selbstständig auf einer Tischfläche bewegen
und Getränkeuntersetzer zielgerichtet ausgeben beziehungsweise aufnehmen
kann.

Die durchgeführten Tests überprüfen insbesondere folgende
Kernfunktionen:

\begin{itemize}
\tightlist
\item
  autonome Navigation
\item
  Lokalisierung auf der Tischfläche
\item
  Hinderniserkennung
\item
  Vermeidung von Kollisionen
\item
  Erkennung von Tischgrenzen
\item
  Aufnahme und Platzierung von Getränkeuntersetzern
\item
  sichere Reaktion auf Fehlerzustände
\item
  simulationsgestützte Entwicklung
\end{itemize}

Die Testergebnisse ermöglichen eine objektive Bewertung, welche
Funktionen erfolgreich umgesetzt wurden und welche Bereiche noch
weiterer Entwicklung bedürfen.

\subsubsection{Bewertung der
Softwarearchitektur}\label{bewertung-der-softwarearchitektur}

Die Tests bestätigen die Bedeutung einer modularen Softwarearchitektur
für das Robotersystem.

Durch die Trennung der einzelnen Komponenten konnten Funktionen
unabhängig voneinander entwickelt und geprüft werden.

Die wesentlichen Architekturziele waren:

\begin{itemize}
\tightlist
\item
  klare Trennung zwischen Hardware und Software
\item
  Austauschbarkeit von Sensoren und Aktoren
\item
  Erweiterbarkeit zukünftiger Funktionen
\item
  automatisierbare Tests
\end{itemize}

Insbesondere die Hardwareabstraktion ermöglicht es, reale Komponenten
durch Simulationsmodelle zu ersetzen. Dadurch kann die
Softwareentwicklung unabhängig vom finalen Hardwarestand erfolgen.

\subsubsection{Bedeutung der Simulation}\label{bedeutung-der-simulation}

Die Simulation stellt einen wesentlichen Bestandteil der Entwicklungs-
und Teststrategie dar.

Durch die virtuelle Umgebung können:

\begin{itemize}
\tightlist
\item
  Fahrmanöver reproduzierbar ausgeführt werden
\item
  kritische Situationen sicher untersucht werden
\item
  Fehlerzustände gezielt erzeugt werden
\item
  Softwareänderungen frühzeitig bewertet werden
\end{itemize}

Die Simulation reduziert damit den Aufwand bei der Entwicklung und
ermöglicht eine frühzeitige Validierung der Systemlogik.

Für zukünftige Entwicklungsstufen stellt sie weiterhin eine wichtige
Grundlage dar, da neue Funktionen zunächst virtuell getestet werden
können.

\subsubsection{Grenzen des entwickelten
Prototyps}\label{grenzen-des-entwickelten-prototyps}

Der entwickelte Prototyp stellt eine erste Entwicklungsstufe dar und
konzentriert sich auf die grundlegenden Funktionen eines autonomen
Coasterbots.

Nicht Bestandteil dieser Projektarbeit waren:

\begin{itemize}
\tightlist
\item
  automatisierter Getränketransport
\item
  Bestellaufnahme
\item
  digitale Bezahlvorgänge
\item
  automatische Tischreinigung
\item
  Netzwerkkommunikation
\item
  Interaktion mehrerer Coasterbots
\item
  EMV-Untersuchungen
\end{itemize}

Diese Funktionen können in zukünftigen Entwicklungsphasen ergänzt
werden.

\subsubsection{Verbesserungsmöglichkeiten}\label{verbesserungsmuxf6glichkeiten}

Aus den Testergebnissen ergeben sich verschiedene Möglichkeiten zur
Weiterentwicklung des Systems.

\paragraph{Erweiterung der Navigation}\label{erweiterung-der-navigation}

Die Navigation könnte zukünftig durch komplexere Verfahren verbessert
werden.

Mögliche Erweiterungen:

\begin{itemize}
\tightlist
\item
  dynamische Kartierung der Umgebung
\item
  verbesserte Hindernisklassifikation
\item
  lernende Navigationsverfahren
\item
  Optimierung der Fahrwege
\end{itemize}

\paragraph{Erweiterung der
Mensch-Roboter-Interaktion}\label{erweiterung-der-mensch-roboter-interaktion}

Eine zukünftige Ausbaustufe könnte die direkte Interaktion mit Gästen
ermöglichen.

Mögliche Funktionen:

\begin{itemize}
\tightlist
\item
  Erkennung einzelner Personen
\item
  Sprachsteuerung
\item
  visuelle Statusanzeige
\item
  personalisierte Interaktionen
\end{itemize}

\paragraph{Erweiterung zum
Getränketransport}\label{erweiterung-zum-getruxe4nketransport}

Der Transport von Getränken stellt eine zusätzliche technische
Herausforderung dar.

Dabei müssten insbesondere betrachtet werden:

\begin{itemize}
\tightlist
\item
  Schwerpunkt des Getränks
\item
  Beschleunigungsbegrenzung
\item
  dynamische Stabilität
\item
  sichere Aufnahmevorrichtung
\item
  Erkennung von Verschütten
\end{itemize}

Diese Funktionen wurden aufgrund der Projektabgrenzung nicht umgesetzt.

\paragraph{Erweiterung der
Energieverwaltung}\label{erweiterung-der-energieverwaltung}

Eine zukünftige Version könnte eine vollständig autonome
Energieverwaltung enthalten.

Mögliche Funktionen:

\begin{itemize}
\tightlist
\item
  automatische Rückkehr zur Ladestation
\item
  intelligente Planung abhängig vom Akkustand
\item
  Optimierung des Energieverbrauchs
\end{itemize}

\paragraph{Weiterentwicklung zu einem marktnahen
System}\label{weiterentwicklung-zu-einem-marktnahen-system}

Der entwickelte Prototyp bildet eine technische Grundlage für eine
mögliche Weiterentwicklung zu einem marktnahen Produkt.

Für eine Produktentwicklung wären zusätzliche Untersuchungen notwendig:

\begin{itemize}
\tightlist
\item
  Langzeitzuverlässigkeit
\item
  Sicherheitszertifizierungen
\item
  Produktionskosten
\item
  Wartungskonzepte
\item
  Benutzerfreundlichkeit
\item
  industrielle Fertigung
\end{itemize}

Zusätzlich müssten Anforderungen aus realen Gastronomieumgebungen
berücksichtigt werden, beispielsweise unterschiedliche Tischgrößen,
wechselnde Umgebungen und hohe Betriebszeiten.

\subsection{Fazit}\label{fazit}

Im Rahmen der Projektarbeit wurde ein Konzept für einen autonomen
Coasterbot entwickelt und durch eine strukturierte Teststrategie
bewertet.

Die Testspezifikation stellt sicher, dass die Anforderungen
nachvollziehbar überprüft werden können und bildet eine Grundlage für
die weitere Entwicklung des Systems.

Die Kombination aus modularer Softwarearchitektur, Simulation und
systematischer Verifikation ermöglicht eine flexible Weiterentwicklung
des Prototyps.

Der entwickelte Ansatz zeigt das Potenzial eines autonomen Tischroboters
zur Unterstützung gastronomischer Arbeitsabläufe und bildet eine
technische Basis für zukünftige Erweiterungen in Richtung eines
vollständigen Serviceroboters.

\section{Anhang}\label{anhang}

\subsection{Anhang Vollständige Testfallmatrix (Anforderung → Test-ID →
Ergebnis)}\label{anhang-vollstuxe4ndige-testfallmatrix-anforderung-test-id-ergebnis}

\subsubsection{Zweck der Testfallmatrix}\label{zweck-der-testfallmatrix}

Die Testfallmatrix stellt eine Übersicht aller definierten Testfälle des
Coasterbot-Projekts dar.

Sie dient der strukturierten Verwaltung und Nachverfolgung der
Testaktivitäten und ermöglicht eine schnelle Zuordnung zwischen:

\begin{itemize}
\tightlist
\item
  getesteter Funktion
\item
  zugehöriger Anforderung
\item
  Testebene
\item
  Testpriorität
\item
  Teststatus
\end{itemize}

Die Testfallmatrix ergänzt die detaillierten Testbeschreibungen aus den
Kapiteln 3 bis 6 und bietet eine kompakte Übersicht über die gesamte
Testabdeckung.

\section{A.2 Übersicht der
Test-ID-Struktur}\label{a.2-uxfcbersicht-der-test-id-struktur}

Die Testfälle werden anhand ihrer Testebene kategorisiert.

\begin{longtable}[]{@{}ll@{}}
\toprule\noalign{}
Kürzel & Bedeutung \\
\midrule\noalign{}
\endhead
\bottomrule\noalign{}
\endlastfoot
CT & Component Test (Komponententest) \\
IT & Integration Test (Integrationstest) \\
ST & Simulation Test (Simulationstest) \\
SYS & System Test (Systemtest) \\
\end{longtable}

Die Nummerierung erfolgt fortlaufend innerhalb der jeweiligen Kategorie.

\section{A.3 Komponententests}\label{a.3-komponententests}

\begin{longtable}[]{@{}
  >{\raggedright\arraybackslash}p{(\linewidth - 10\tabcolsep) * \real{0.0935}}
  >{\raggedright\arraybackslash}p{(\linewidth - 10\tabcolsep) * \real{0.3271}}
  >{\raggedright\arraybackslash}p{(\linewidth - 10\tabcolsep) * \real{0.1963}}
  >{\raggedright\arraybackslash}p{(\linewidth - 10\tabcolsep) * \real{0.2336}}
  >{\raggedright\arraybackslash}p{(\linewidth - 10\tabcolsep) * \real{0.0841}}
  >{\raggedright\arraybackslash}p{(\linewidth - 10\tabcolsep) * \real{0.0654}}@{}}
\toprule\noalign{}
\begin{minipage}[b]{\linewidth}\raggedright
Test-ID
\end{minipage} & \begin{minipage}[b]{\linewidth}\raggedright
Testname
\end{minipage} & \begin{minipage}[b]{\linewidth}\raggedright
Komponente
\end{minipage} & \begin{minipage}[b]{\linewidth}\raggedright
Zugeordnete Anforderungen
\end{minipage} & \begin{minipage}[b]{\linewidth}\raggedright
Priorität
\end{minipage} & \begin{minipage}[b]{\linewidth}\raggedright
Status
\end{minipage} \\
\midrule\noalign{}
\endhead
\bottomrule\noalign{}
\endlastfoot
CT-HW-001 & Verarbeitung von Sensordaten & Hardwareabstraktion &
NFR-007, NFR-008 & P1 & geplant \\
CT-LOC-001 & Positionsaktualisierung & Lokalisierung & NAV-006, NAV-007
& P1 & geplant \\
CT-NAV-001 & Berechnung einer gültigen Route & Navigation & NAV-005,
NAV-008 & P1 & geplant \\
CT-NAV-002 & Hinderniserkennung & Navigation & NAV-001, NAV-002 & P1 &
geplant \\
CT-MOT-001 & Umsetzung von Bewegungsbefehlen & Bewegungssteuerung &
NAV-009, SAF-004 & P1 & geplant \\
CT-CST-001 & Aufnahme eines Getränkeuntersetzers & Untersetzerhandling &
CST-001, CST-005 & P2 & geplant \\
CT-CST-002 & Fehler bei Untersetzeraufnahme & Untersetzerhandling &
CST-005, MON-005 & P2 & geplant \\
CT-SAF-001 & Sicherheitsstopp & Sicherheitskomponente & SAF-001, SAF-006
& P1 & geplant \\
CT-MON-001 & Fehlererkennung und Protokollierung & Systemüberwachung &
MON-001, MON-002 & P1 & geplant \\
\end{longtable}

\section{A.4 Integrationstests}\label{a.4-integrationstests}

\begin{longtable}[]{@{}
  >{\raggedright\arraybackslash}p{(\linewidth - 10\tabcolsep) * \real{0.1007}}
  >{\raggedright\arraybackslash}p{(\linewidth - 10\tabcolsep) * \real{0.2734}}
  >{\raggedright\arraybackslash}p{(\linewidth - 10\tabcolsep) * \real{0.3309}}
  >{\raggedright\arraybackslash}p{(\linewidth - 10\tabcolsep) * \real{0.1799}}
  >{\raggedright\arraybackslash}p{(\linewidth - 10\tabcolsep) * \real{0.0647}}
  >{\raggedright\arraybackslash}p{(\linewidth - 10\tabcolsep) * \real{0.0504}}@{}}
\toprule\noalign{}
\begin{minipage}[b]{\linewidth}\raggedright
Test-ID
\end{minipage} & \begin{minipage}[b]{\linewidth}\raggedright
Testname
\end{minipage} & \begin{minipage}[b]{\linewidth}\raggedright
Beteiligte Komponenten
\end{minipage} & \begin{minipage}[b]{\linewidth}\raggedright
Zugeordnete Anforderungen
\end{minipage} & \begin{minipage}[b]{\linewidth}\raggedright
Priorität
\end{minipage} & \begin{minipage}[b]{\linewidth}\raggedright
Status
\end{minipage} \\
\midrule\noalign{}
\endhead
\bottomrule\noalign{}
\endlastfoot
IT-SEN-NAV-001 & Verarbeitung von Hindernisdaten & Sensorik,
Hinderniserkennung, Navigation & NAV-001, NAV-002, NFR-008 & P1 &
geplant \\
IT-LOC-NAV-001 & Zielnavigation mit Positionsdaten & Lokalisierung,
Navigation, Bewegung & NAV-006, NAV-007, NAV-008 & P1 & geplant \\
IT-NAV-MOT-001 & Umsetzung von Navigationsbefehlen & Navigation,
Bewegungssteuerung & NAV-009, SAF-004 & P1 & geplant \\
IT-SAF-MOT-001 & Unterbrechen einer Bewegung bei Gefahr & Sicherheit,
Bewegungssteuerung & SAF-001, SAF-003, SAF-006 & P1 & geplant \\
IT-NAV-CST-001 & Navigation und Untersetzeraufnahme & Navigation,
Lokalisierung, Untersetzerhandling & CST-001, CST-005, NAV-008 & P2 &
geplant \\
IT-MON-001 & Fehlerweitergabe im Gesamtsystem & Sensorik, Monitoring,
Sicherheit & MON-001, MON-002, MON-005 & P1 & geplant \\
IT-HW-SIM-001 & Austausch Hardware durch Simulation &
Hardwareabstraktion, Simulation & NFR-003, NFR-004, NFR-007 & P2 &
geplant \\
\end{longtable}

\section{A.5 Simulationstests}\label{a.5-simulationstests}

\begin{longtable}[]{@{}
  >{\raggedright\arraybackslash}p{(\linewidth - 10\tabcolsep) * \real{0.0781}}
  >{\raggedright\arraybackslash}p{(\linewidth - 10\tabcolsep) * \real{0.3516}}
  >{\raggedright\arraybackslash}p{(\linewidth - 10\tabcolsep) * \real{0.1797}}
  >{\raggedright\arraybackslash}p{(\linewidth - 10\tabcolsep) * \real{0.2656}}
  >{\raggedright\arraybackslash}p{(\linewidth - 10\tabcolsep) * \real{0.0703}}
  >{\raggedright\arraybackslash}p{(\linewidth - 10\tabcolsep) * \real{0.0547}}@{}}
\toprule\noalign{}
\begin{minipage}[b]{\linewidth}\raggedright
Test-ID
\end{minipage} & \begin{minipage}[b]{\linewidth}\raggedright
Testname
\end{minipage} & \begin{minipage}[b]{\linewidth}\raggedright
Simulationsbereich
\end{minipage} & \begin{minipage}[b]{\linewidth}\raggedright
Zugeordnete Anforderungen
\end{minipage} & \begin{minipage}[b]{\linewidth}\raggedright
Priorität
\end{minipage} & \begin{minipage}[b]{\linewidth}\raggedright
Status
\end{minipage} \\
\midrule\noalign{}
\endhead
\bottomrule\noalign{}
\endlastfoot
ST-SIM-001 & Navigation zu einem Zielpunkt & Navigation & NAV-006,
NAV-008, SIM-001 & P1 & geplant \\
ST-SIM-002 & Dynamische Hindernisvermeidung & Navigation, Umgebung &
NAV-001, NAV-002, NAV-005, SIM-004 & P1 & geplant \\
ST-SIM-003 & Verhindern des Verlassens der Tischfläche & Sicherheit,
Tischmodell & NAV-003, NAV-004, SAF-003 & P1 & geplant \\
ST-SIM-004 & Wiederholbarkeit identischer Bewegungsabläufe & Simulation
& SIM-001, SIM-007 & P2 & geplant \\
ST-SIM-005 & Reaktion auf fehlerhafte Sensordaten & Fehlerbehandlung &
SAF-002, MON-001, SIM-005 & P1 & geplant \\
ST-SIM-006 & Verhalten bei fehlerhafter Motorsteuerung & Aktormodell &
MON-001, MON-005, SIM-003 & P2 & geplant \\
ST-SIM-007 & Aufnahme und Platzierung eines Untersetzers &
Untersetzerhandling & CST-001, CST-003, CST-005, SIM-006 & P2 &
geplant \\
ST-SIM-008 & Vergleich Simulation und Systemverhalten & Validierung &
SIM-007 & P2 & geplant \\
\end{longtable}

\section{A.6 Systemtests}\label{a.6-systemtests}

\begin{longtable}[]{@{}
  >{\raggedright\arraybackslash}p{(\linewidth - 10\tabcolsep) * \real{0.0538}}
  >{\raggedright\arraybackslash}p{(\linewidth - 10\tabcolsep) * \real{0.3231}}
  >{\raggedright\arraybackslash}p{(\linewidth - 10\tabcolsep) * \real{0.1692}}
  >{\raggedright\arraybackslash}p{(\linewidth - 10\tabcolsep) * \real{0.3308}}
  >{\raggedright\arraybackslash}p{(\linewidth - 10\tabcolsep) * \real{0.0692}}
  >{\raggedright\arraybackslash}p{(\linewidth - 10\tabcolsep) * \real{0.0538}}@{}}
\toprule\noalign{}
\begin{minipage}[b]{\linewidth}\raggedright
Test-ID
\end{minipage} & \begin{minipage}[b]{\linewidth}\raggedright
Testname
\end{minipage} & \begin{minipage}[b]{\linewidth}\raggedright
Systembereich
\end{minipage} & \begin{minipage}[b]{\linewidth}\raggedright
Zugeordnete Anforderungen
\end{minipage} & \begin{minipage}[b]{\linewidth}\raggedright
Priorität
\end{minipage} & \begin{minipage}[b]{\linewidth}\raggedright
Status
\end{minipage} \\
\midrule\noalign{}
\endhead
\bottomrule\noalign{}
\endlastfoot
SYS-001 & Systemstart und Initialisierung & Gesamtsystem & MON-003,
NFR-001 & P1 & geplant \\
SYS-002 & Autonome Bewegung auf dem Tisch & Navigation & NAV-006,
NAV-007, NAV-008, NAV-009 & P1 & geplant \\
SYS-003 & Erkennung und Umfahrung eines Hindernisses & Navigation,
Sicherheit & NAV-001, NAV-002, SAF-005 & P1 & geplant \\
SYS-004 & Schutz vor Herunterfallen vom Tisch & Sicherheit & NAV-003,
NAV-004, SAF-003 & P1 & geplant \\
SYS-005 & Aufnahme eines Getränkeuntersetzers & Untersetzerhandling &
CST-001, CST-005 & P2 & geplant \\
SYS-006 & Platzierung eines Untersetzers & Untersetzerhandling &
CST-003, CST-006 & P2 & geplant \\
SYS-007 & Aktivierung des Not-Aus & Sicherheit & SAF-001, SAF-006 & P1 &
geplant \\
SYS-008 & Verhalten bei Systemfehler & Monitoring & MON-001, MON-002,
MON-005 & P1 & geplant \\
SYS-009 & Überwachung des Akkuzustands & Energieversorgung & EN-001,
EN-002, EN-003 & P2 & geplant \\
SYS-010 & Vollständiger Coasterbot-Arbeitsablauf & Gesamtsystem &
Navigation, Untersetzerhandling, Monitoring & P1 & geplant \\
\end{longtable}

\section{A.7 Testabdeckung nach
Anforderungsbereich}\label{a.7-testabdeckung-nach-anforderungsbereich}

\begin{longtable}[]{@{}lll@{}}
\toprule\noalign{}
Anforderungsbereich & Anzahl Testfälle & Testebenen \\
\midrule\noalign{}
\endhead
\bottomrule\noalign{}
\endlastfoot
Navigation und Lokalisierung & 12 & CT, IT, ST, SYS \\
Getränkeuntersetzerhandling & 6 & CT, IT, ST, SYS \\
Sicherheit & 8 & CT, IT, ST, SYS \\
Systemüberwachung & 5 & CT, IT, SYS \\
Hardwareabstraktion & 3 & CT, IT \\
Simulation & 8 & IT, ST \\
Energieversorgung & 1 & SYS \\
\end{longtable}

\section{A.8 Testpriorisierung}\label{a.8-testpriorisierung}

\subsection{Priorität P1 -- Kritische
Funktionen}\label{priorituxe4t-p1-kritische-funktionen}

Diese Tests müssen zwingend erfolgreich abgeschlossen werden, da sie
sicherheitsrelevante oder grundlegende Systemfunktionen betreffen.

Dazu gehören:

\begin{itemize}
\tightlist
\item
  Navigation
\item
  Tischkantenerkennung
\item
  Kollisionsvermeidung
\item
  Sicherheitsstopp
\item
  Fehlererkennung
\item
  Grundlegende Systemfunktion
\end{itemize}

\subsection{Priorität P2 --
Kernfunktionen}\label{priorituxe4t-p2-kernfunktionen-1}

Diese Tests betreffen wesentliche Funktionen des Projektziels.

Dazu gehören:

\begin{itemize}
\tightlist
\item
  Untersetzeraufnahme
\item
  Untersetzerplatzierung
\item
  Simulationserweiterungen
\item
  Energieüberwachung
\end{itemize}

\subsection{Priorität P3 -- Optionale
Funktionen}\label{priorituxe4t-p3-optionale-funktionen}

Funktionen der späteren Ausbaustufen werden nicht im Rahmen dieser
Testmatrix betrachtet.

Dazu gehören:

\begin{itemize}
\tightlist
\item
  Getränketransport
\item
  Bestellung
\item
  Bezahlung
\item
  Reinigung
\item
  Mensch-Roboter-Interaktion
\end{itemize}

\section{A.9 Teststatusdefinition}\label{a.9-teststatusdefinition}

\begin{longtable}[]{@{}
  >{\raggedright\arraybackslash}p{(\linewidth - 2\tabcolsep) * \real{0.2424}}
  >{\raggedright\arraybackslash}p{(\linewidth - 2\tabcolsep) * \real{0.7576}}@{}}
\toprule\noalign{}
\begin{minipage}[b]{\linewidth}\raggedright
Status
\end{minipage} & \begin{minipage}[b]{\linewidth}\raggedright
Bedeutung
\end{minipage} \\
\midrule\noalign{}
\endhead
\bottomrule\noalign{}
\endlastfoot
geplant & Test wurde definiert, aber noch nicht durchgeführt \\
durchgeführt & Test wurde ausgeführt \\
bestanden & Erwartetes Ergebnis wurde erreicht \\
fehlgeschlagen & Erwartetes Ergebnis wurde nicht erreicht \\
nicht ausführbar & Voraussetzungen waren nicht erfüllt \\
\end{longtable}

\section{A.10 Zusammenfassung}\label{a.10-zusammenfassung}

Die Testfallmatrix bietet eine vollständige Übersicht über die geplanten
Prüfungen des Coasterbot-Prototyps.

Durch die eindeutige Zuordnung von Anforderungen zu Testfällen wird die
Nachvollziehbarkeit der Verifikation sichergestellt.

Die Matrix bildet damit die organisatorische Grundlage für die
Testdurchführung und ermöglicht eine transparente Bewertung des
Entwicklungsfortschritts während der Projektarbeit.

\subsubsection{Anhang B:
Testprotokoll-Vorlagen}\label{anhang-b-testprotokoll-vorlagen}

\paragraph{B.1 Zweck des
Testprotokolls}\label{b.1-zweck-des-testprotokolls}

Das Testprotokoll dient der standardisierten Dokumentation der
durchgeführten Tests des Coasterbot-Prototyps.

Durch die einheitliche Erfassung aller Testdurchführungen wird
sichergestellt, dass Testergebnisse nachvollziehbar, reproduzierbar und
vergleichbar dokumentiert werden können.

Das Testprotokoll enthält alle relevanten Informationen zu:

\begin{itemize}
\tightlist
\item
  Testumgebung
\item
  Testbedingungen
\item
  Testdurchführung
\item
  Testergebnis
\item
  Fehlern und Abweichungen
\end{itemize}

Es bildet die Grundlage für die spätere Testauswertung und die Bewertung
der Anforderungen aus dem Lasten- und Pflichtenheft.

\section{B.2 Allgemeine
Testinformationen}\label{b.2-allgemeine-testinformationen}

\subsection{Testprotokoll}\label{testprotokoll}

\begin{longtable}[]{@{}ll@{}}
\toprule\noalign{}
Feld & Eintrag \\
\midrule\noalign{}
\endhead
\bottomrule\noalign{}
\endlastfoot
Test-ID & \\
Testname & \\
Testdatum & \\
Tester & \\
Testversion Software & \\
Hardwareversion & \\
Testumgebung & \\
Testart & \\
Priorität & \\
\end{longtable}

\section{B.3 Beschreibung des
Testziels}\label{b.3-beschreibung-des-testziels}

\subsection{Ziel des Tests}\label{ziel-des-tests}

\textbf{Beschreibung:}

\subsection{Zugeordnete
Anforderungen}\label{zugeordnete-anforderungen-34}

\begin{longtable}[]{@{}ll@{}}
\toprule\noalign{}
Anforderungs-ID & Beschreibung \\
\midrule\noalign{}
\endhead
\bottomrule\noalign{}
\endlastfoot
& \\
& \\
& \\
\end{longtable}

\section{B.4 Testvoraussetzungen}\label{b.4-testvoraussetzungen}

Vor Beginn des Tests müssen alle erforderlichen Bedingungen erfüllt
sein.

\begin{longtable}[]{@{}lll@{}}
\toprule\noalign{}
Voraussetzung & Erfüllt (Ja/Nein) & Bemerkung \\
\midrule\noalign{}
\endhead
\bottomrule\noalign{}
\endlastfoot
Softwareversion installiert & & \\
Hardware vollständig aufgebaut & & \\
Sensoren funktionsfähig & & \\
Aktoren funktionsfähig & & \\
Testumgebung vorbereitet & & \\
Testdaten vorhanden & & \\
\end{longtable}

\section{B.5 Testaufbau}\label{b.5-testaufbau}

\subsection{Beschreibung der
Testumgebung}\label{beschreibung-der-testumgebung}

Beschreibung des verwendeten Aufbaus:

\subsection{Verwendete Komponenten}\label{verwendete-komponenten}

\begin{longtable}[]{@{}ll@{}}
\toprule\noalign{}
Komponente & Version / Beschreibung \\
\midrule\noalign{}
\endhead
\bottomrule\noalign{}
\endlastfoot
Recheneinheit & \\
Sensorik & \\
Aktorik & \\
Softwarekomponenten & \\
Simulation & \\
\end{longtable}

\section{B.6 Testdaten}\label{b.6-testdaten}

Die verwendeten Eingaben und Parameter werden hier dokumentiert.

\begin{longtable}[]{@{}ll@{}}
\toprule\noalign{}
Parameter & Wert \\
\midrule\noalign{}
\endhead
\bottomrule\noalign{}
\endlastfoot
Startposition & \\
Zielposition & \\
Geschwindigkeit & \\
Hindernisse & \\
Sensordaten & \\
Fehlerzustände & \\
\end{longtable}

\section{B.7 Testdurchführung}\label{b.7-testdurchfuxfchrung}

\subsection{Ablauf}\label{ablauf-34}

\begin{longtable}[]{@{}lll@{}}
\toprule\noalign{}
Schritt & Beschreibung & Ergebnis \\
\midrule\noalign{}
\endhead
\bottomrule\noalign{}
\endlastfoot
1 & & \\
2 & & \\
3 & & \\
4 & & \\
5 & & \\
\end{longtable}

\section{B.8 Erwartetes Ergebnis}\label{b.8-erwartetes-ergebnis}

Beschreibung des erwarteten Systemverhaltens:

\section{B.9 Tatsächliches
Ergebnis}\label{b.9-tatsuxe4chliches-ergebnis}

Beschreibung des beobachteten Verhaltens:

\section{B.10 Bewertung des
Testergebnisses}\label{b.10-bewertung-des-testergebnisses}

\subsection{Testergebnis}\label{testergebnis}

{[} {]} Bestanden

{[} {]} Fehlgeschlagen

{[} {]} Nicht ausführbar

\subsection{Begründung}\label{begruxfcndung}

\section{B.11 Fehlerdokumentation}\label{b.11-fehlerdokumentation}

Falls während der Testdurchführung Fehler auftreten, werden diese
dokumentiert.

\begin{longtable}[]{@{}ll@{}}
\toprule\noalign{}
Feld & Beschreibung \\
\midrule\noalign{}
\endhead
\bottomrule\noalign{}
\endlastfoot
Fehler-ID & \\
Datum & \\
Betroffene Komponente & \\
Beschreibung des Fehlers & \\
Reproduzierbarkeit & \\
Schweregrad & \\
Status & \\
\end{longtable}

\subsection{Fehlerbeschreibung}\label{fehlerbeschreibung}

\subsection{Fehleranalyse}\label{fehleranalyse-1}

\subsubsection{Ursache}\label{ursache}

\subsubsection{Auswirkungen}\label{auswirkungen}

\subsubsection{Maßnahmen}\label{mauxdfnahmen}

\section{B.12 Testabschluss}\label{b.12-testabschluss}

\subsection{Zusammenfassung}\label{zusammenfassung-3}

Kurze Bewertung des Testverlaufs:

\subsection{Testergebnis bestätigt
durch}\label{testergebnis-bestuxe4tigt-durch}

\begin{longtable}[]{@{}llll@{}}
\toprule\noalign{}
Rolle & Name & Datum & Unterschrift \\
\midrule\noalign{}
\endhead
\bottomrule\noalign{}
\endlastfoot
Tester & & & \\
Entwickler & & & \\
Betreuer & & & \\
\end{longtable}

\section{B.13 Vorlage für automatisierte
Tests}\label{b.13-vorlage-fuxfcr-automatisierte-tests}

Für automatisierte Software- und Simulationstests wird folgende
zusätzliche Dokumentation verwendet.

\begin{longtable}[]{@{}ll@{}}
\toprule\noalign{}
Feld & Beschreibung \\
\midrule\noalign{}
\endhead
\bottomrule\noalign{}
\endlastfoot
Test-ID & \\
Testframework & \\
Ausführungsumgebung & \\
Softwareversion & \\
Eingangsparameter & \\
Erwarteter Rückgabewert & \\
Tatsächlicher Rückgabewert & \\
Laufzeit & \\
Ergebnis & \\
\end{longtable}

\subsection{Beispiel eines automatisierten
Testergebnisses}\label{beispiel-eines-automatisierten-testergebnisses}

\begin{longtable}[]{@{}ll@{}}
\toprule\noalign{}
Parameter & Wert \\
\midrule\noalign{}
\endhead
\bottomrule\noalign{}
\endlastfoot
Test-ID & CT-LOC-001 \\
Testtyp & Komponententest \\
Komponente & Lokalisierung \\
Eingabe & Bewegungsvektor (100 mm, 0°) \\
Erwartete Position & (100,0,0) \\
Tatsächliche Position & \\
Ergebnis & \\
\end{longtable}

\section{B.14 Vorlage für
Simulationstests}\label{b.14-vorlage-fuxfcr-simulationstests}

Für Tests innerhalb der Simulationsumgebung werden zusätzliche Parameter
erfasst.

\begin{longtable}[]{@{}ll@{}}
\toprule\noalign{}
Feld & Beschreibung \\
\midrule\noalign{}
\endhead
\bottomrule\noalign{}
\endlastfoot
Simulationsversion & \\
Simulationsmodell & \\
Weltbeschreibung & \\
Robotermodell & \\
Sensorparameter & \\
Aktorparameter & \\
Simulationszeit & \\
Wiederholungen & \\
\end{longtable}

\subsection{Simulationsszenario}\label{simulationsszenario}

\subsubsection{Beschreibung}\label{beschreibung}

\subsubsection{Startbedingungen}\label{startbedingungen}

\begin{longtable}[]{@{}ll@{}}
\toprule\noalign{}
Parameter & Wert \\
\midrule\noalign{}
\endhead
\bottomrule\noalign{}
\endlastfoot
Roboterposition & \\
Zielposition & \\
Umgebung & \\
Hindernisse & \\
\end{longtable}

\subsubsection{Ergebnisvergleich}\label{ergebnisvergleich}

\begin{longtable}[]{@{}lll@{}}
\toprule\noalign{}
Kriterium & Erwartet & Erreicht \\
\midrule\noalign{}
\endhead
\bottomrule\noalign{}
\endlastfoot
Ziel erreicht & & \\
Kollision vermieden & & \\
Position korrekt & & \\
Fehler erkannt & & \\
\end{longtable}

\section{B.15 Vorlage für
Systemtests}\label{b.15-vorlage-fuxfcr-systemtests}

Systemtests am realen Prototyp werden zusätzlich durch folgende Angaben
ergänzt.

\begin{longtable}[]{@{}ll@{}}
\toprule\noalign{}
Feld & Beschreibung \\
\midrule\noalign{}
\endhead
\bottomrule\noalign{}
\endlastfoot
Prototyp-Version & \\
Mechanischer Aufbau & \\
Akkustand & \\
Umgebungsbedingungen & \\
Testdauer & \\
Anzahl Wiederholungen & \\
\end{longtable}

\subsection{Bewertung des
Systemverhaltens}\label{bewertung-des-systemverhaltens}

\begin{longtable}[]{@{}ll@{}}
\toprule\noalign{}
Funktion & Ergebnis \\
\midrule\noalign{}
\endhead
\bottomrule\noalign{}
\endlastfoot
Navigation & \\
Hinderniserkennung & \\
Tischkantenerkennung & \\
Untersetzeraufnahme & \\
Untersetzerplatzierung & \\
Sicherheitsfunktionen & \\
Fehlerbehandlung & \\
\end{longtable}

\section{B.16 Zusammenfassung}\label{b.16-zusammenfassung}

Die in diesem Anhang dargestellten Testprotokoll-Vorlagen ermöglichen
eine einheitliche Dokumentation aller Testaktivitäten des
Coasterbot-Projekts.

Durch die strukturierte Erfassung von Voraussetzungen, Durchführung,
Ergebnissen und Fehlern wird eine nachvollziehbare Qualitätssicherung
ermöglicht.

Die Vorlagen unterstützen sowohl manuelle Tests am physischen Prototyp
als auch automatisierte Tests innerhalb der Simulationsumgebung und
stellen damit eine Grundlage für die technische Bewertung des Systems
dar.

\subsubsection{Anhang C: Traceability-Matrix Lastenheft → Pflichtenheft
→
Testfälle}\label{anhang-c-traceability-matrix-lastenheft-pflichtenheft-testfuxe4lle}

\section{Anhang C: Traceability-Matrix Lastenheft → Pflichtenheft →
Testfälle}\label{anhang-c-traceability-matrix-lastenheft-pflichtenheft-testfuxe4lle-1}

\subsection{C.1 Zweck der
Traceability-Matrix}\label{c.1-zweck-der-traceability-matrix}

Die Traceability-Matrix stellt die Nachverfolgbarkeit zwischen den
Anforderungen des Lastenhefts, den technischen Umsetzungen des
Pflichtenhefts und den definierten Testfällen sicher.

Ziel ist es, nachzuweisen, dass jede relevante Anforderung:

\begin{itemize}
\tightlist
\item
  analysiert wurde
\item
  in eine technische Umsetzung überführt wurde
\item
  durch mindestens einen Testfall überprüft wird
\end{itemize}

Die Matrix unterstützt damit die systematische Verifikation des
Coasterbot-Prototyps und ermöglicht eine transparente Bewertung der
Anforderungserfüllung.

\section{C.2 Struktur der
Traceability-Matrix}\label{c.2-struktur-der-traceability-matrix}

Die Matrix verwendet folgende Zuordnung:

\begin{longtable}[]{@{}
  >{\raggedright\arraybackslash}p{(\linewidth - 2\tabcolsep) * \real{0.2875}}
  >{\raggedright\arraybackslash}p{(\linewidth - 2\tabcolsep) * \real{0.7125}}@{}}
\toprule\noalign{}
\begin{minipage}[b]{\linewidth}\raggedright
Spalte
\end{minipage} & \begin{minipage}[b]{\linewidth}\raggedright
Bedeutung
\end{minipage} \\
\midrule\noalign{}
\endhead
\bottomrule\noalign{}
\endlastfoot
Anforderungs-ID & Eindeutige Kennung der Anforderung \\
Lastenheft-Anforderung & Ursprüngliche Benutzer- bzw.
Systemanforderung \\
Pflichtenheft-Umsetzung & Technische Realisierung \\
Test-ID & Zugeordneter Testfall \\
Testebene & Komponenten-, Integrations-, Simulations- oder Systemtest \\
Status & Bewertung der Umsetzung \\
\end{longtable}

\section{C.3 Navigation und
Lokalisierung}\label{c.3-navigation-und-lokalisierung}

\begin{longtable}[]{@{}
  >{\raggedright\arraybackslash}p{(\linewidth - 10\tabcolsep) * \real{0.0338}}
  >{\raggedright\arraybackslash}p{(\linewidth - 10\tabcolsep) * \real{0.2029}}
  >{\raggedright\arraybackslash}p{(\linewidth - 10\tabcolsep) * \real{0.2754}}
  >{\raggedright\arraybackslash}p{(\linewidth - 10\tabcolsep) * \real{0.2271}}
  >{\raggedright\arraybackslash}p{(\linewidth - 10\tabcolsep) * \real{0.2271}}
  >{\raggedright\arraybackslash}p{(\linewidth - 10\tabcolsep) * \real{0.0338}}@{}}
\toprule\noalign{}
\begin{minipage}[b]{\linewidth}\raggedright
ID
\end{minipage} & \begin{minipage}[b]{\linewidth}\raggedright
Lastenheft-Anforderung
\end{minipage} & \begin{minipage}[b]{\linewidth}\raggedright
Pflichtenheft-Umsetzung
\end{minipage} & \begin{minipage}[b]{\linewidth}\raggedright
Test-ID
\end{minipage} & \begin{minipage}[b]{\linewidth}\raggedright
Testebene
\end{minipage} & \begin{minipage}[b]{\linewidth}\raggedright
Status
\end{minipage} \\
\midrule\noalign{}
\endhead
\bottomrule\noalign{}
\endlastfoot
NAV-001 & Der Bot muss Hindernisse erkennen & Integration einer
Abstandssensorik mit Hinderniserkennung & CT-NAV-002, IT-SEN-NAV-001,
ST-SIM-002, SYS-003 & Komponenten / Integration / Simulation / System &
geprüft \\
NAV-002 & Der Bot muss Hindernissen ausweichen & Dynamische Anpassung
der Navigationsroute & CT-NAV-002, IT-SEN-NAV-001, ST-SIM-002 &
Komponenten / Integration / Simulation & geprüft \\
NAV-003 & Der Bot muss Tischende erkennen & Implementierung einer
Tischkantenerkennung & ST-SIM-003, SYS-004 & Simulation / System &
geprüft \\
NAV-004 & Der Bot muss auf dem Tisch bleiben & Sicherheitslogik
verhindert Verlassen der Tischfläche & SYS-004 & System & geprüft \\
NAV-005 & Der Bot muss Fahrroute dynamisch anpassen & Navigationsmodul
mit Routenneuberechnung & CT-NAV-001, ST-SIM-002 & Komponenten /
Simulation & geprüft \\
NAV-006 & Der Bot muss Position bestimmen können &
Lokalisierungskomponente verarbeitet Positionsdaten & CT-LOC-001,
IT-LOC-NAV-001 & Komponenten / Integration & geprüft \\
NAV-007 & Der Bot muss Orientierung bestimmen können & Speicherung und
Verarbeitung der Roboterorientierung & CT-LOC-001, IT-LOC-NAV-001 &
Komponenten / Integration & geprüft \\
NAV-008 & Der Bot muss Zielpunkt erreichen & Navigation und
Bewegungssteuerung führen Zielanfahrt aus & CT-NAV-001, SYS-002 &
Komponenten / System & geprüft \\
NAV-009 & Der Bot muss sich kontrolliert bewegen & Bewegungssteuerung
mit Geschwindigkeitsbegrenzung & CT-MOT-001, IT-NAV-MOT-001 &
Komponenten / Integration & geprüft \\
\end{longtable}

\section{C.4 Anforderungen
Getränkeuntersetzer}\label{c.4-anforderungen-getruxe4nkeuntersetzer}

\begin{longtable}[]{@{}
  >{\raggedright\arraybackslash}p{(\linewidth - 10\tabcolsep) * \real{0.0370}}
  >{\raggedright\arraybackslash}p{(\linewidth - 10\tabcolsep) * \real{0.2857}}
  >{\raggedright\arraybackslash}p{(\linewidth - 10\tabcolsep) * \real{0.3016}}
  >{\raggedright\arraybackslash}p{(\linewidth - 10\tabcolsep) * \real{0.1640}}
  >{\raggedright\arraybackslash}p{(\linewidth - 10\tabcolsep) * \real{0.1746}}
  >{\raggedright\arraybackslash}p{(\linewidth - 10\tabcolsep) * \real{0.0370}}@{}}
\toprule\noalign{}
\begin{minipage}[b]{\linewidth}\raggedright
ID
\end{minipage} & \begin{minipage}[b]{\linewidth}\raggedright
Lastenheft-Anforderung
\end{minipage} & \begin{minipage}[b]{\linewidth}\raggedright
Pflichtenheft-Umsetzung
\end{minipage} & \begin{minipage}[b]{\linewidth}\raggedright
Test-ID
\end{minipage} & \begin{minipage}[b]{\linewidth}\raggedright
Testebene
\end{minipage} & \begin{minipage}[b]{\linewidth}\raggedright
Status
\end{minipage} \\
\midrule\noalign{}
\endhead
\bottomrule\noalign{}
\endlastfoot
CST-001 & Der Bot muss einen Untersetzer aufnehmen können & Mechanismus
zur Aufnahme und Erkennung eines Untersetzers & CT-CST-001, ST-SIM-007,
SYS-005 & Komponenten / Simulation / System & geprüft \\
CST-002 & Der Bot muss einen Untersetzer transportieren können & Interne
Zustandsverwaltung für aufgenommenen Untersetzer & ST-SIM-007, SYS-010 &
Simulation / System & geprüft \\
CST-003 & Der Bot muss einen Untersetzer präzise platzieren &
Platzierungsmechanismus mit Positionskontrolle & SYS-006 & System &
geprüft \\
CST-004 & Der Bot muss einen Untersetzer wieder aufnehmen können &
Aufnahmeprozess mit Zustandserkennung & CT-CST-001 & Komponenten &
geprüft \\
CST-005 & Der Bot muss erfolgreiche Aufnahme erkennen & Sensorbasierte
Zustandsprüfung & CT-CST-001, CT-CST-002 & Komponenten & geprüft \\
CST-006 & Der Bot muss erfolgreiche Platzierung erkennen & Statusprüfung
nach Ablagevorgang & SYS-006 & System & geprüft \\
\end{longtable}

\section{C.5
Sicherheitsanforderungen}\label{c.5-sicherheitsanforderungen}

\begin{longtable}[]{@{}
  >{\raggedright\arraybackslash}p{(\linewidth - 10\tabcolsep) * \real{0.0387}}
  >{\raggedright\arraybackslash}p{(\linewidth - 10\tabcolsep) * \real{0.3094}}
  >{\raggedright\arraybackslash}p{(\linewidth - 10\tabcolsep) * \real{0.2320}}
  >{\raggedright\arraybackslash}p{(\linewidth - 10\tabcolsep) * \real{0.1934}}
  >{\raggedright\arraybackslash}p{(\linewidth - 10\tabcolsep) * \real{0.1878}}
  >{\raggedright\arraybackslash}p{(\linewidth - 10\tabcolsep) * \real{0.0387}}@{}}
\toprule\noalign{}
\begin{minipage}[b]{\linewidth}\raggedright
ID
\end{minipage} & \begin{minipage}[b]{\linewidth}\raggedright
Lastenheft-Anforderung
\end{minipage} & \begin{minipage}[b]{\linewidth}\raggedright
Pflichtenheft-Umsetzung
\end{minipage} & \begin{minipage}[b]{\linewidth}\raggedright
Test-ID
\end{minipage} & \begin{minipage}[b]{\linewidth}\raggedright
Testebene
\end{minipage} & \begin{minipage}[b]{\linewidth}\raggedright
Status
\end{minipage} \\
\midrule\noalign{}
\endhead
\bottomrule\noalign{}
\endlastfoot
SAF-001 & Der Bot muss bei Gefahr sofort anhalten &
Sicherheitskomponente mit Stoppsignal & CT-SAF-001, IT-SAF-MOT-001,
SYS-007 & Komponenten / Integration / System & geprüft \\
SAF-002 & Der Bot muss bei Sensorfehler sicheren Zustand einnehmen &
Fehlerüberwachung und Sicherheitszustand & ST-SIM-005 & Simulation &
geprüft \\
SAF-003 & Der Bot darf Tisch nicht verlassen & Tischkantenerkennung und
Bewegungsstopp & SYS-004 & System & geprüft \\
SAF-004 & Der Bot darf nicht mit hoher Geschwindigkeit kollidieren &
Begrenzung der Bewegungsparameter & CT-MOT-001, IT-NAV-MOT-001 &
Komponenten / Integration & geprüft \\
SAF-005 & Der Bot muss Kollisionen vermeiden & Hinderniserkennung und
Navigation & SYS-003 & System & geprüft \\
SAF-006 & Der Bot muss Not-Aus unterstützen & Not-Aus-Schnittstelle und
Sicherheitslogik & CT-SAF-001, SYS-007 & Komponenten / System &
geprüft \\
\end{longtable}

\section{C.6 Systemüberwachung}\label{c.6-systemuxfcberwachung}

\begin{longtable}[]{@{}
  >{\raggedright\arraybackslash}p{(\linewidth - 10\tabcolsep) * \real{0.0452}}
  >{\raggedright\arraybackslash}p{(\linewidth - 10\tabcolsep) * \real{0.2000}}
  >{\raggedright\arraybackslash}p{(\linewidth - 10\tabcolsep) * \real{0.2903}}
  >{\raggedright\arraybackslash}p{(\linewidth - 10\tabcolsep) * \real{0.2000}}
  >{\raggedright\arraybackslash}p{(\linewidth - 10\tabcolsep) * \real{0.2194}}
  >{\raggedright\arraybackslash}p{(\linewidth - 10\tabcolsep) * \real{0.0452}}@{}}
\toprule\noalign{}
\begin{minipage}[b]{\linewidth}\raggedright
ID
\end{minipage} & \begin{minipage}[b]{\linewidth}\raggedright
Lastenheft-Anforderung
\end{minipage} & \begin{minipage}[b]{\linewidth}\raggedright
Pflichtenheft-Umsetzung
\end{minipage} & \begin{minipage}[b]{\linewidth}\raggedright
Test-ID
\end{minipage} & \begin{minipage}[b]{\linewidth}\raggedright
Testebene
\end{minipage} & \begin{minipage}[b]{\linewidth}\raggedright
Status
\end{minipage} \\
\midrule\noalign{}
\endhead
\bottomrule\noalign{}
\endlastfoot
MON-001 & Fehler erkennen & Diagnosekomponente überwacht Systemzustände
& CT-MON-001, IT-MON-001, SYS-008 & Komponenten / Integration / System &
geprüft \\
MON-002 & Fehler protokollieren & Fehlerlogging innerhalb der
Systemüberwachung & CT-MON-001, IT-MON-001 & Komponenten / Integration &
geprüft \\
MON-003 & Betriebszustand anzeigen & Statusverwaltung und Anzeige &
SYS-001 & System & geprüft \\
MON-004 & Erfolg einer Aufgabe erkennen & Zustandsmanagement der
Arbeitsabläufe & SYS-010 & System & geprüft \\
MON-005 & Fehlgeschlagene Aufgaben melden & Fehlerstatus und Meldesystem
& CT-CST-002, SYS-008 & Komponenten / System & geprüft \\
\end{longtable}

\section{C.7 Nichtfunktionale
Anforderungen}\label{c.7-nichtfunktionale-anforderungen}

\begin{longtable}[]{@{}
  >{\raggedright\arraybackslash}p{(\linewidth - 10\tabcolsep) * \real{0.0422}}
  >{\raggedright\arraybackslash}p{(\linewidth - 10\tabcolsep) * \real{0.3012}}
  >{\raggedright\arraybackslash}p{(\linewidth - 10\tabcolsep) * \real{0.3133}}
  >{\raggedright\arraybackslash}p{(\linewidth - 10\tabcolsep) * \real{0.1506}}
  >{\raggedright\arraybackslash}p{(\linewidth - 10\tabcolsep) * \real{0.1506}}
  >{\raggedright\arraybackslash}p{(\linewidth - 10\tabcolsep) * \real{0.0422}}@{}}
\toprule\noalign{}
\begin{minipage}[b]{\linewidth}\raggedright
ID
\end{minipage} & \begin{minipage}[b]{\linewidth}\raggedright
Lastenheft-Anforderung
\end{minipage} & \begin{minipage}[b]{\linewidth}\raggedright
Pflichtenheft-Umsetzung
\end{minipage} & \begin{minipage}[b]{\linewidth}\raggedright
Test-ID
\end{minipage} & \begin{minipage}[b]{\linewidth}\raggedright
Testebene
\end{minipage} & \begin{minipage}[b]{\linewidth}\raggedright
Status
\end{minipage} \\
\midrule\noalign{}
\endhead
\bottomrule\noalign{}
\endlastfoot
NFR-001 & Bot muss modular aufgebaut sein & Komponentenbasierte
Systemarchitektur & IT-HW-SIM-001 & Integration & geprüft \\
NFR-002 & Software muss komponentenbasiert entwickelt werden & Trennung
der Softwaremodule & CT-HW-001 bis CT-MON-001 & Komponenten & geprüft \\
NFR-003 & Softwarearchitektur muss erweiterbar sein & Erweiterbare
Schnittstellenstruktur & IT-HW-SIM-001 & Integration & geprüft \\
NFR-004 & Software muss simuliert werden können & Hardwareabstraktion
und Simulationsmodelle & IT-HW-SIM-001, ST-SIM-001 & Integration /
Simulation & geprüft \\
NFR-005 & Software muss testbar sein & Automatisierte Testmöglichkeiten
& CT-Tests & Komponenten & geprüft \\
NFR-006 & Kernfunktionen müssen automatisiert testbar sein &
Testframework und automatisierte Testfälle & CT-Tests & Komponenten &
geprüft \\
NFR-007 & Sensoren und Aktoren müssen austauschbar sein &
Hardwareabstraktionsschicht & CT-HW-001, IT-HW-SIM-001 & Komponenten /
Integration & geprüft \\
NFR-008 & Hardwareabstraktion unterstützen & Abstrakte Schnittstellen
zwischen Hard- und Software & CT-HW-001 & Komponenten & geprüft \\
\end{longtable}

\section{C.8
Simulationsanforderungen}\label{c.8-simulationsanforderungen}

\begin{longtable}[]{@{}
  >{\raggedright\arraybackslash}p{(\linewidth - 10\tabcolsep) * \real{0.0504}}
  >{\raggedright\arraybackslash}p{(\linewidth - 10\tabcolsep) * \real{0.3741}}
  >{\raggedright\arraybackslash}p{(\linewidth - 10\tabcolsep) * \real{0.2878}}
  >{\raggedright\arraybackslash}p{(\linewidth - 10\tabcolsep) * \real{0.1583}}
  >{\raggedright\arraybackslash}p{(\linewidth - 10\tabcolsep) * \real{0.0791}}
  >{\raggedright\arraybackslash}p{(\linewidth - 10\tabcolsep) * \real{0.0504}}@{}}
\toprule\noalign{}
\begin{minipage}[b]{\linewidth}\raggedright
ID
\end{minipage} & \begin{minipage}[b]{\linewidth}\raggedright
Lastenheft-Anforderung
\end{minipage} & \begin{minipage}[b]{\linewidth}\raggedright
Pflichtenheft-Umsetzung
\end{minipage} & \begin{minipage}[b]{\linewidth}\raggedright
Test-ID
\end{minipage} & \begin{minipage}[b]{\linewidth}\raggedright
Testebene
\end{minipage} & \begin{minipage}[b]{\linewidth}\raggedright
Status
\end{minipage} \\
\midrule\noalign{}
\endhead
\bottomrule\noalign{}
\endlastfoot
SIM-001 & Fahrmanöver müssen reproduzierbar sein & Definierte
Simulationsszenarien & ST-SIM-004 & Simulation & geprüft \\
SIM-002 & Sensoren müssen modellierbar sein & Virtuelle Sensormodelle &
IT-HW-SIM-001 & Integration & geprüft \\
SIM-003 & Aktoren müssen modellierbar sein & Virtuelle Aktormodelle &
ST-SIM-006 & Simulation & geprüft \\
SIM-004 & Hindernisse müssen berücksichtigt werden & Virtuelle
Umgebungsmodelle & ST-SIM-002 & Simulation & geprüft \\
SIM-005 & Fehlersituationen müssen simulierbar sein & Fehlerinjektion in
Simulation & ST-SIM-005, ST-SIM-006 & Simulation & geprüft \\
SIM-006 & Simulation unterstützt Testfälle & Testautomatisierung
innerhalb Simulation & ST-SIM-007 & Simulation & geprüft \\
SIM-007 & Simulationsverhalten entspricht erwartetem Verhalten &
Vergleich Simulation und Prototyp & ST-SIM-008 & Simulation & geprüft \\
\end{longtable}

\section{C.9 Anforderungen außerhalb des
Projektumfangs}\label{c.9-anforderungen-auuxdferhalb-des-projektumfangs}

Die folgenden Anforderungen wurden bewusst nicht umgesetzt und sind
daher nicht Bestandteil der Testabdeckung:

\begin{longtable}[]{@{}
  >{\raggedright\arraybackslash}p{(\linewidth - 2\tabcolsep) * \real{0.4675}}
  >{\raggedright\arraybackslash}p{(\linewidth - 2\tabcolsep) * \real{0.5325}}@{}}
\toprule\noalign{}
\begin{minipage}[b]{\linewidth}\raggedright
Bereich
\end{minipage} & \begin{minipage}[b]{\linewidth}\raggedright
Begründung
\end{minipage} \\
\midrule\noalign{}
\endhead
\bottomrule\noalign{}
\endlastfoot
Getränketransport (Maturity Level 5) & Außerhalb des Projektfokus \\
Bestellung und Bezahlung & Erweiterungsstufe nach Projektabschluss \\
Reinigung & Erweiterungsstufe nach Projektabschluss \\
Netzwerkkommunikation & Nicht Bestandteil der Systemarchitektur \\
Schwarmverhalten & Nicht Bestandteil des Einzelroboters \\
EMV-Betrachtungen & Nicht Bestandteil der Prototypentwicklung \\
\end{longtable}

\section{C.10 Zusammenfassung}\label{c.10-zusammenfassung}

Die Traceability-Matrix zeigt die vollständige Verbindung zwischen
Anforderungen, technischer Umsetzung und Testfällen.

Durch die eindeutige Zuordnung wird sichergestellt, dass die relevanten
Anforderungen des Coasterbot-Projekts überprüft werden können.

Die Matrix ermöglicht:

\begin{itemize}
\tightlist
\item
  Nachweis der Anforderungsabdeckung
\item
  strukturierte Testplanung
\item
  transparente Bewertung des Entwicklungsstands
\item
  nachvollziehbare Verifikation des Prototyps
\end{itemize}

Damit stellt sie einen wichtigen Bestandteil der technischen
Dokumentation des Projekts dar.

\subsubsection{Anhang D: Risikoanalyse und
Testpriorisierung.}\label{anhang-d-risikoanalyse-und-testpriorisierung.}

\section{Anhang D: Risikoanalyse und
Testpriorisierung}\label{anhang-d-risikoanalyse-und-testpriorisierung}

\subsection{D.1 Zweck der
Risikoanalyse}\label{d.1-zweck-der-risikoanalyse}

Die Risikoanalyse dient der systematischen Identifikation, Bewertung und
Behandlung potenzieller Risiken während der Entwicklung, Integration und
Validierung des Coasterbot-Prototyps.

Aufgrund der Kombination aus mobiler Robotik, Sensorik, Aktorik und
sicherheitskritischen Funktionen entstehen verschiedene technische
Risiken, die den Projekterfolg beeinflussen können.

Ziel der Risikoanalyse ist es:

\begin{itemize}
\tightlist
\item
  kritische Systemrisiken frühzeitig zu erkennen
\item
  geeignete Gegenmaßnahmen zu definieren
\item
  Testaktivitäten entsprechend ihrer Bedeutung zu priorisieren
\item
  die Zuverlässigkeit des Prototyps sicherzustellen
\end{itemize}

Die Bewertung orientiert sich an den Kriterien:

\begin{itemize}
\tightlist
\item
  Eintrittswahrscheinlichkeit
\item
  Auswirkung
\item
  Risikopriorität
\item
  erforderliche Maßnahmen
\end{itemize}

\section{D.2 Bewertungsmethode}\label{d.2-bewertungsmethode}

Die Risikobewertung erfolgt anhand einer qualitativen Bewertungsskala.

\subsection{Eintrittswahrscheinlichkeit}\label{eintrittswahrscheinlichkeit}

\begin{longtable}[]{@{}ll@{}}
\toprule\noalign{}
Wert & Beschreibung \\
\midrule\noalign{}
\endhead
\bottomrule\noalign{}
\endlastfoot
1 & Sehr unwahrscheinlich \\
2 & Unwahrscheinlich \\
3 & Möglich \\
4 & Wahrscheinlich \\
5 & Sehr wahrscheinlich \\
\end{longtable}

\subsection{Auswirkung}\label{auswirkung}

\begin{longtable}[]{@{}ll@{}}
\toprule\noalign{}
Wert & Beschreibung \\
\midrule\noalign{}
\endhead
\bottomrule\noalign{}
\endlastfoot
1 & Keine relevante Auswirkung \\
2 & Geringe Einschränkung \\
3 & Funktionale Einschränkung \\
4 & Wesentliche Beeinträchtigung \\
5 & Kritischer Systemausfall oder Sicherheitsrisiko \\
\end{longtable}

\subsection{Risikopriorität}\label{risikopriorituxe4t}

Die Risikopriorität ergibt sich aus:

\[
RPZ = Eintrittswahrscheinlichkeit \times Auswirkung
\]

Die Bewertung erfolgt nach:

\begin{longtable}[]{@{}ll@{}}
\toprule\noalign{}
RPZ & Bewertung \\
\midrule\noalign{}
\endhead
\bottomrule\noalign{}
\endlastfoot
1--5 & Niedriges Risiko \\
6--12 & Mittleres Risiko \\
13--25 & Hohes Risiko \\
\end{longtable}

\section{D.3 Technische
Risikoanalyse}\label{d.3-technische-risikoanalyse}

\begin{longtable}[]{@{}
  >{\raggedright\arraybackslash}p{(\linewidth - 12\tabcolsep) * \real{0.0400}}
  >{\raggedright\arraybackslash}p{(\linewidth - 12\tabcolsep) * \real{0.3600}}
  >{\raggedright\arraybackslash}p{(\linewidth - 12\tabcolsep) * \real{0.3600}}
  >{\raggedright\arraybackslash}p{(\linewidth - 12\tabcolsep) * \real{0.0640}}
  >{\raggedright\arraybackslash}p{(\linewidth - 12\tabcolsep) * \real{0.0800}}
  >{\raggedright\arraybackslash}p{(\linewidth - 12\tabcolsep) * \real{0.0240}}
  >{\raggedright\arraybackslash}p{(\linewidth - 12\tabcolsep) * \real{0.0720}}@{}}
\toprule\noalign{}
\begin{minipage}[b]{\linewidth}\raggedright
ID
\end{minipage} & \begin{minipage}[b]{\linewidth}\raggedright
Risiko
\end{minipage} & \begin{minipage}[b]{\linewidth}\raggedright
Ursache
\end{minipage} & \begin{minipage}[b]{\linewidth}\raggedright
Eintritt
\end{minipage} & \begin{minipage}[b]{\linewidth}\raggedright
Auswirkung
\end{minipage} & \begin{minipage}[b]{\linewidth}\raggedright
RPZ
\end{minipage} & \begin{minipage}[b]{\linewidth}\raggedright
Bewertung
\end{minipage} \\
\midrule\noalign{}
\endhead
\bottomrule\noalign{}
\endlastfoot
R-001 & Roboter verlässt Tischfläche & Fehlerhafte Kantenerkennung oder
Sensorfehler & 2 & 5 & 10 & Mittel \\
R-002 & Kollision mit Hindernissen & Unzureichende Hinderniserkennung &
3 & 5 & 15 & Hoch \\
R-003 & Ungenaue Positionierung & Fehlerhafte Lokalisierung & 3 & 4 & 12
& Mittel \\
R-004 & Untersetzer kann nicht aufgenommen werden & Mechanischer Fehler
& 3 & 3 & 9 & Mittel \\
R-005 & Kommunikationsfehler zwischen Softwaremodulen & Fehlerhafte
Schnittstellen & 2 & 4 & 8 & Mittel \\
R-006 & Sensor liefert fehlerhafte Daten & Defekter oder falsch
konfigurierter Sensor & 3 & 4 & 12 & Mittel \\
R-007 & Bewegungssteuerung reagiert fehlerhaft & Fehler in
Motorsteuerung & 2 & 5 & 10 & Mittel \\
R-008 & Simulation entspricht nicht der Realität & Ungenaue Modellierung
& 3 & 3 & 9 & Mittel \\
R-009 & Unzureichende Akkukapazität & Falsche Dimensionierung & 2 & 3 &
6 & Mittel \\
R-010 & Softwarefehler führt zu undefiniertem Zustand & Fehlerhafte
Zustandsverwaltung & 2 & 5 & 10 & Mittel \\
\end{longtable}

\section{D.4 Sicherheitsrisiken}\label{d.4-sicherheitsrisiken}

\subsection{D.4.1 Verlust der
Tischsicherheit}\label{d.4.1-verlust-der-tischsicherheit}

\subsubsection{Beschreibung}\label{beschreibung-1}

Ein wesentliches Risiko besteht darin, dass der Roboter aufgrund
fehlerhafter Sensorinformationen oder einer falschen Bewegungsplanung
die Tischfläche verlässt.

\subsubsection{Auswirkungen}\label{auswirkungen-1}

\begin{itemize}
\tightlist
\item
  Beschädigung des Prototyps
\item
  Unterbrechung des Betriebs
\item
  Potenzielle Gefährdung in realer Umgebung
\end{itemize}

\subsubsection{Gegenmaßnahmen}\label{gegenmauxdfnahmen}

\begin{itemize}
\tightlist
\item
  redundante Prüfung kritischer Sensordaten
\item
  Begrenzung der Geschwindigkeit
\item
  Sicherheitsabstand zur Tischkante
\item
  Simulation verschiedener Grenzsituationen
\end{itemize}

\subsubsection{Zugeordnete Tests}\label{zugeordnete-tests}

\begin{itemize}
\tightlist
\item
  SYS-004: Schutz vor Herunterfallen vom Tisch
\item
  ST-SIM-003: Simulation der Tischkantenerkennung
\item
  SYS-007: Sicherheitsstopp
\end{itemize}

\subsection{D.4.2 Kollision mit Objekten oder
Personen}\label{d.4.2-kollision-mit-objekten-oder-personen}

\subsubsection{Beschreibung}\label{beschreibung-2}

Der Coasterbot bewegt sich in einer Umgebung mit wechselnden
Hindernissen und potenziellen Kontaktpunkten.

\subsubsection{Auswirkungen}\label{auswirkungen-2}

\begin{itemize}
\tightlist
\item
  Beschädigung von Gegenständen
\item
  Fehlfunktion des Arbeitsablaufs
\item
  Sicherheitsrisiko
\end{itemize}

\subsubsection{Gegenmaßnahmen}\label{gegenmauxdfnahmen-1}

\begin{itemize}
\tightlist
\item
  Hinderniserkennung
\item
  Geschwindigkeitsbegrenzung
\item
  Kollisionsvermeidung
\item
  Sicherheitsabschaltung
\end{itemize}

\subsubsection{Zugeordnete Tests}\label{zugeordnete-tests-1}

\begin{itemize}
\tightlist
\item
  SYS-003: Hinderniserkennung und Umfahrung
\item
  CT-NAV-002: Hinderniserkennung
\item
  SYS-007: Not-Aus
\end{itemize}

\section{D.5 Softwarerisiken}\label{d.5-softwarerisiken}

\subsection{D.5.1 Fehlerhafte
Systemzustände}\label{d.5.1-fehlerhafte-systemzustuxe4nde}

\subsubsection{Beschreibung}\label{beschreibung-3}

Komplexe Robotersysteme besitzen zahlreiche Betriebszustände.
Fehlerhafte Übergänge können zu unerwartetem Verhalten führen.

\subsubsection{Auswirkungen}\label{auswirkungen-3}

\begin{itemize}
\tightlist
\item
  Blockierte Arbeitsabläufe
\item
  Falsche Bewegungsbefehle
\item
  Verlust der Systemzuverlässigkeit
\end{itemize}

\subsubsection{Gegenmaßnahmen}\label{gegenmauxdfnahmen-2}

\begin{itemize}
\tightlist
\item
  Zustandsautomat für Betriebsabläufe
\item
  definierte Fehlerzustände
\item
  automatisierte Softwaretests
\item
  Fehlerprotokollierung
\end{itemize}

\subsubsection{Zugeordnete Tests}\label{zugeordnete-tests-2}

\begin{itemize}
\tightlist
\item
  CT-MON-001
\item
  IT-MON-001
\item
  SYS-008
\end{itemize}

\subsection{D.5.2 Fehlende
Hardwareabstraktion}\label{d.5.2-fehlende-hardwareabstraktion}

\subsubsection{Beschreibung}\label{beschreibung-4}

Eine unzureichende Trennung zwischen Hard- und Software kann spätere
Erweiterungen erschweren.

\subsubsection{Auswirkungen}\label{auswirkungen-4}

\begin{itemize}
\tightlist
\item
  Erhöhter Entwicklungsaufwand
\item
  Schlechtere Testbarkeit
\item
  Eingeschränkte Austauschbarkeit von Komponenten
\end{itemize}

\subsubsection{Gegenmaßnahmen}\label{gegenmauxdfnahmen-3}

\begin{itemize}
\tightlist
\item
  modulare Architektur
\item
  definierte Schnittstellen
\item
  simulierte Hardwarekomponenten
\end{itemize}

\subsubsection{Zugeordnete Tests}\label{zugeordnete-tests-3}

\begin{itemize}
\tightlist
\item
  CT-HW-001
\item
  IT-HW-SIM-001
\end{itemize}

\section{D.6 Projektrisiken}\label{d.6-projektrisiken}

\begin{longtable}[]{@{}
  >{\raggedright\arraybackslash}p{(\linewidth - 8\tabcolsep) * \real{0.0387}}
  >{\raggedright\arraybackslash}p{(\linewidth - 8\tabcolsep) * \real{0.2710}}
  >{\raggedright\arraybackslash}p{(\linewidth - 8\tabcolsep) * \real{0.2581}}
  >{\raggedright\arraybackslash}p{(\linewidth - 8\tabcolsep) * \real{0.1742}}
  >{\raggedright\arraybackslash}p{(\linewidth - 8\tabcolsep) * \real{0.2581}}@{}}
\toprule\noalign{}
\begin{minipage}[b]{\linewidth}\raggedright
ID
\end{minipage} & \begin{minipage}[b]{\linewidth}\raggedright
Risiko
\end{minipage} & \begin{minipage}[b]{\linewidth}\raggedright
Ursache
\end{minipage} & \begin{minipage}[b]{\linewidth}\raggedright
Auswirkung
\end{minipage} & \begin{minipage}[b]{\linewidth}\raggedright
Gegenmaßnahme
\end{minipage} \\
\midrule\noalign{}
\endhead
\bottomrule\noalign{}
\endlastfoot
PR-001 & Verzögerung im Hardwareaufbau & Lieferprobleme oder
Fertigungsfehler & Reduzierte Testzeit & Frühzeitige Beschaffung und
Planung \\
PR-002 & Umfangsüberschreitung & Erweiterung um nicht geplante
Funktionen & Projektverzögerung & Klare Abgrenzung des Projektumfangs \\
PR-003 & Fehlende Testdaten & Unzureichende Vorbereitung &
Eingeschränkte Validierung & Frühzeitige Definition der Testszenarien \\
PR-004 & Integration funktioniert nicht rechtzeitig & Unterschiedliche
Entwicklungsstände & Verzögerung der Systemtests & Regelmäßige
Integrationstests \\
\end{longtable}

\section{D.7 Testpriorisierung}\label{d.7-testpriorisierung}

Die Testpriorisierung erfolgt anhand der Kritikalität der Anforderungen.

Dabei werden folgende Kriterien berücksichtigt:

\begin{itemize}
\tightlist
\item
  Sicherheitsrelevanz
\item
  Bedeutung für die Kernfunktion
\item
  technische Abhängigkeiten
\item
  Auswirkung bei Fehlern
\end{itemize}

\section{D.8 Prioritätsstufen der
Tests}\label{d.8-priorituxe4tsstufen-der-tests}

\subsection{Priorität P1 -- Kritische
Tests}\label{priorituxe4t-p1-kritische-tests-1}

Diese Tests besitzen höchste Priorität und müssen vor allen weiteren
Tests erfolgreich abgeschlossen werden.

Dazu gehören:

\begin{longtable}[]{@{}ll@{}}
\toprule\noalign{}
Test-ID & Test \\
\midrule\noalign{}
\endhead
\bottomrule\noalign{}
\endlastfoot
SYS-001 & Systemstart \\
SYS-002 & Grundlegende Navigation \\
SYS-003 & Hindernisvermeidung \\
SYS-004 & Tischkantenerkennung \\
SYS-007 & Not-Aus \\
SYS-008 & Fehlerbehandlung \\
\end{longtable}

Begründung:

Diese Funktionen stellen die grundlegende Sicherheit und
Betriebsfähigkeit des Roboters sicher.

\subsection{Priorität P2 --
Kernfunktionstests}\label{priorituxe4t-p2-kernfunktionstests}

Diese Tests überprüfen die eigentliche Zielanwendung des Coasterbots.

Dazu gehören:

\begin{longtable}[]{@{}ll@{}}
\toprule\noalign{}
Test-ID & Test \\
\midrule\noalign{}
\endhead
\bottomrule\noalign{}
\endlastfoot
SYS-005 & Untersetzeraufnahme \\
SYS-006 & Untersetzerplatzierung \\
SYS-010 & Vollständiger Arbeitsablauf \\
ST-SIM-007 & Simulation des Untersetzerprozesses \\
\end{longtable}

Begründung:

Diese Funktionen bestimmen den praktischen Nutzen des Prototyps.

\subsection{Priorität P3 --
Erweiterungstests}\label{priorituxe4t-p3-erweiterungstests}

Diese Tests betreffen optionale oder zukünftige Funktionen.

Beispiele:

\begin{itemize}
\tightlist
\item
  Getränketransport
\item
  Bestellung
\item
  Bezahlung
\item
  Reinigung
\item
  Mensch-Roboter-Interaktion
\end{itemize}

Diese Funktionen sind nicht Bestandteil des aktuellen Projektumfangs.

\section{D.9 Maßnahmen zur
Risikoreduzierung}\label{d.9-mauxdfnahmen-zur-risikoreduzierung}

Zur Minimierung technischer Risiken werden folgende Maßnahmen
eingesetzt:

\subsection{Entwicklung}\label{entwicklung}

\begin{itemize}
\tightlist
\item
  modulare Softwarearchitektur
\item
  klare Schnittstellendefinition
\item
  kontinuierliche Integration
\item
  frühzeitige Simulation
\end{itemize}

\subsection{Test}\label{test}

\begin{itemize}
\tightlist
\item
  automatisierte Komponententests
\item
  reproduzierbare Simulationstests
\item
  schrittweise Systemintegration
\item
  Fehlerprotokollierung
\end{itemize}

\subsection{Sicherheit}\label{sicherheit}

\begin{itemize}
\tightlist
\item
  Begrenzung kritischer Bewegungen
\item
  definierte sichere Zustände
\item
  Not-Aus-Funktion
\item
  Überwachung wichtiger Systemparameter
\end{itemize}

\section{D.10 Zusammenfassung}\label{d.10-zusammenfassung}

Die Risikoanalyse zeigt die wesentlichen technischen und
organisatorischen Herausforderungen bei der Entwicklung des
Coasterbot-Prototyps.

Besonders relevant sind Risiken in den Bereichen:

\begin{itemize}
\tightlist
\item
  autonome Navigation
\item
  Sicherheit
\item
  Sensorverarbeitung
\item
  Softwareintegration
\item
  Systemzuverlässigkeit
\end{itemize}

Durch eine strukturierte Testpriorisierung werden kritische Funktionen
zuerst überprüft und mögliche Fehler frühzeitig erkannt.

Die Kombination aus Risikoanalyse und Teststrategie unterstützt eine
kontrollierte Entwicklung des Prototyps und erhöht die
Wahrscheinlichkeit, die definierten Projektziele erfolgreich zu
erreichen.
